% 农业上市公司画像：动物保健Ⅲ、动物保健（共 9 家，按流通市值降序）
% 由 scripts/make_company_profiles.py 自动生成，请勿手工修改

\subsubsection{生物股份（600201）}
\label{co:600201}

\pfl{公司基本情况} 内蒙古呼和浩特的兽用生物制品企业，全称金宇生物技术股份有限公司，1992 年 12 月设立、1999 年 1 月 15 日在上海证券交易所上市；主营兽用生物制品的研发、生产与销售，产品涵盖猪、禽、反刍和宠物类四大系列百余种动物疫苗，拥有口蹄疫、高致病性禽流感和布鲁氏菌病三大强制免疫疫苗农业农村部定点生产资质。2025 年归母净利润 1.00 亿元、同比下降 8.14\%；2026 年上半年归母净利润 1.04 亿元、同比增长 72.28\%。公司两次再融资（2021 年度、2023 年度）均告终止。 公司于 1999-01-15 在上交所首发上市，注册资本 11.00 亿元，总股本 11.20 亿股。 截至 2026 年 9 月 24 日收盘价 11.40 元，流通市值 125.4 亿元，近一年+28.2\%，流通市值在三级行业“动物保健Ⅲ”的 9 家公司中按由大到小排第\zhnumber{1}位，近一年涨跌幅高于全样本中位数（-10.1\%）38.3 个百分点。 按 2025 年年报披露的行业构成，收入占比前三位为生物制药 92.7\%；其他(补充) 3.9\%；其他收入 3.5\%。 股权结构方面，控股股东为内蒙古金宇生物控股有限公司。

\pfl{历史股价} 自 2021 年 1 月至 2026 年 9 月，股价由 22.77 元变为 11.40 元，累计 -49.9\%，区间最高 22.77 元（2021 年 1 月）、最低 6.26 元（2024 年 8 月）；当前价相当于区间最高价的 50.1\%，为区间最低价的 1.82 倍。 月度收益的年化波动率 37.5\%，高于全样本中位数 37.4\%，在三级行业“动物保健Ⅲ”的 9 家公司中波动率由高到低排第\zhnumber{5}位。 把股价阶段与公司事件对齐看，2024 年 8 月 至 2026 年 1 月股价上涨+188.3\%，同期（2025 年 7 月）原拟向控股股东方非公开发行不超过 8.00 亿元（发行价 8.29 元/股）的计划已于 2025 年 7 月终止；2021 年 1 月 至 2022 年 9 月股价下跌-65.7\%，同期（2022 年 4 月）公司终止 2021 年度非公开发行 A 股股票事项。 从时间对应关系看，公司层面的资本运作与股权、风险事件解释了区间内多数急涨急跌，与行业指数的同步性相对有限。

\begin{center}
\begin{minipage}{\textwidth}\centering
\begin{tikzpicture}
\begin{axis}[
  width=12.6cm, height=3.9cm, scale only axis,
  xmin=2020.922, xmax=2026.888, ymin=5.82, ymax=24.36,
  xtick={2021,2022,2023,2024,2025,2026},
  yearticks,
  yticklabel style={font=\scriptsize},
  ylabel={元/股}, ylabel style={font=\scriptsize},
  grid=major, grid style={LightGray, line width=0.3pt},
  axis line style={InkGray, line width=0.5pt},
  every axis plot/.append style={line width=0.9pt},
]
\addplot[AgriGreen] table[x=x,y=y]{data/clean/profiles/px_600201.dat};
\addplot[SoftRed, dashed, line width=0.7pt] table[x=x,y=y]{data/clean/profiles/pxzz_600201.dat};
\addplot[only marks, mark=*, mark size=1.1pt, SoftRed] table[x=x,y=y]{data/clean/profiles/pxzz_600201.dat};
\end{axis}
\end{tikzpicture}

\captionof{figure}[生物股份（600201）股价与周期骨架]{生物股份（600201）月度收盘价（元/股）与 ZigZag 周期骨架（阈值 25\%，圆点为周期转折点）}
\end{minipage}
\end{center}

\pfl{过去周期分析} 以 25\% 的 ZigZag 阈值划分，样本区间内股价形成 2 个主升段与 3 个主跌段。 最大主升段出现在 2024 年 8 月 至 2026 年 1 月，17 个月+188.3\%；最大主跌段为 2021 年 1 月 至 2022 年 9 月，20 个月-65.7\%。 最近一次转折出现在 2026 年 9 月（下跌段自 2026 年 1 月起，8 个月-36.8\%），当前处于该段的延续之中。 公司股价较申万农林牧渔指数提前约 21 个月见底，是板块中较早定价周期反转的公司之一。 动保的需求由存栏量而非价格决定（第 \ref{sec:vet} 节）：猪价下行阶段龙头养殖企业的逆势扩产反而推升了疫苗与兽药需求，这正是动保股价与猪价“脱钩”的原因，公司 2026 年上半年实现盈利、半年 ROE 1.9\%（未年化），好于全样本中位数（0.75\%），下行周期对其报表的冲击小于同业。 综合区间形态看，该股单段涨跌幅度偏大、以趋势段为主（最大主升 +188.3\%、最大主跌 -65.7\%），年化波动率 37.5\% 与全样本中位数 37.4\% 相比波动与样本中位数接近。

\pfl{盈利情况} 2026 年上半年营业收入 7.56 亿元（同比 +22.1\%），归母净利润 1.04 亿元，同比 +72.3\%。 以年报为趋势对照，2023---2025 年营业收入由 15.98 亿元变为 13.96 亿元（两年年均复合增速 -6.5\%），归母净利润由 2.84 亿元变为 1.00 亿元。 按半年报口径，2026 年上半年的 ROE 为 1.9\%（未年化）、销售净利率 14.7\%、毛利率 53.9\%，ROE 在“动物保健Ⅲ”9 家公司中按数值由高到低排第\zhnumber{7}位。 同期全样本半年 ROE 中位数 0.75\%，公司与之相差 1.17 个百分点（略好于中位数公司）。 综合看，公司在报告期维持盈利（高于全样本半年 ROE 中位数 0.75\%），利润的可持续性取决于动物保健Ⅲ景气的延续与成本管控的执行。

\begin{center}\small\setlength{\tabcolsep}{3.4pt}
\captionof{table}[生物股份（600201）关键财务指标]{生物股份（600201）关键财务指标（合并报表口径；两个半年报期均未年化；现金短债比 $=$ 货币资金 $\div$ 短期债务）}
\begin{adjustbox}{max width=\textwidth}
\begin{tabular}{@{}lrrrrr@{}}
\toprule
指标 & 2023 年 & 2024 年 & 2025 年 & 2025 年半年报 & 2026 年半年报 \\
\midrule
营业收入（亿元） & 15.98 & 12.55 & 13.96 & 6.20 & 7.56 \\
归母净利润（亿元） & 2.84 & 1.09 & 1.00 & 0.60 & 1.04 \\
毛利率（\%） & 59.2 & 53.9 & 53.5 & 52.2 & 53.9 \\
ROE（\%） & 5.4 & 1.9 & 1.9 & 1.1 & 1.9 \\
资产负债率（\%） & 17.7 & 19.9 & 16.0 & 16.8 & 16.5 \\
经营活动现金流净额（亿元） & 4.15 & 2.90 & 3.16 & 0.89 & 0.51 \\
货币资金（亿元） & 17.12 & 11.40 & 12.65 & 10.60 & 16.09 \\
有息负债合计（亿元） & 0.04 & 0.18 & 0.02 & 0.17 & 1.01 \\
现金短债比（倍） & 1,117.89 & 69.10 & 884.75 & 64.56 & 1,426.53 \\
\bottomrule
\end{tabular}
\end{adjustbox}
\end{center}

\begin{center}
\begin{minipage}{\textwidth}\centering
\begin{tikzpicture}
\begin{groupplot}[
  group style={group size=2 by 1, horizontal sep=0.60cm},
  width=6.85cm, height=4.60cm,
  tick label style={font=\tiny},
  title style={font=\scriptsize\heiti},
  yticklabel style={font=\tiny},
  ylabel style={font=\tiny},
  grid=major, grid style={LightGray, line width=0.3pt},
  axis line style={InkGray, line width=0.5pt},
  enlarge x limits={abs=0.8},
  legend style={font=\scriptsize, at={(0.5,-0.20)}, anchor=north,
                legend columns=2, draw=none, fill=none, column sep=6pt},
]
\nextgroupplot[
  ybar, bar width=4pt,
  ymin=-1.78, ymax=19.89,
  xtick={0,1,2,3,4,5.7},
  xticklabels={2021,2022,2023,2024,2025,2026H1},
  ylabel={亿元},
]
\addplot[fill=AgriGreen!75, draw=AgriGreen, bar shift=-2.6pt] table[x=i, y=rev]{data/clean/profiles/fin_600201.dat};
\addplot[fill=SoftRed!60, draw=SoftRed, bar shift=2.6pt] table[x=i, y=ni]{data/clean/profiles/fin_600201.dat};
\addplot[fill=AgriGreen!30, draw=AgriGreen, pattern=north east lines, bar shift=-2.6pt] table[x=x, y=rev]{data/clean/profiles/finh_600201.dat};
\addplot[fill=SoftRed!20, draw=SoftRed, pattern=north east lines, bar shift=2.6pt] table[x=x, y=ni]{data/clean/profiles/finh_600201.dat};
\legend{营业收入（年/半年）, 归母净利润（年/半年）}
\nextgroupplot[
  ymin=-1.6, ymax=21.9,
  xtick={0,1,2,3,4,5.7},
  xticklabels={2021,2022,2023,2024,2025,2026H1},
  ylabel={\%},
]
\addplot[AgriGold, mark=*, mark size=1.3pt] table[x=i, y=roe]{data/clean/profiles/fin_600201.dat};
\addplot[OilBlue, mark=square*, mark size=1.3pt] table[x=i, y=debt]{data/clean/profiles/fin_600201.dat};
\addplot[only marks, AgriGold, mark=*, mark size=2.0pt] table[x=x, y=roe]{data/clean/profiles/finh_600201.dat};
\addplot[only marks, OilBlue, mark=square*, mark size=2.0pt] table[x=x, y=debt]{data/clean/profiles/finh_600201.dat};
\legend{ROE, 资产负债率}
\end{groupplot}
\end{tikzpicture}
\begin{tikzpicture}
\begin{axis}[
  name=finbot,
  width=13.75cm, height=3.10cm, scale only axis,
  ymin=-1.71, ymax=19.18,
  xtick={0,1,2,3,4,5.7},
  xticklabels={2021,2022,2023,2024,2025,2026H1},
  tick label style={font=\tiny},
  yticklabel style={font=\tiny},
  ylabel={亿元}, ylabel style={font=\tiny},
  ybar, bar width=4pt,
  grid=major, grid style={LightGray, line width=0.3pt},
  axis line style={InkGray, line width=0.5pt},
  enlarge x limits={abs=0.8},
  legend style={font=\scriptsize, at={(0.5,-0.26)}, anchor=north,
                legend columns=2, draw=none, fill=none, column sep=10pt},
]
\addplot[fill=AgriGreen!55, draw=AgriGreen, bar shift=-3.0pt] table[x=i, y=cash]{data/clean/profiles/fin_600201.dat};
\addplot[fill=SoftRed!45, draw=SoftRed, bar shift=3.0pt] table[x=i, y=ibd]{data/clean/profiles/fin_600201.dat};
\addplot[fill=AgriGreen!25, draw=AgriGreen, pattern=north east lines, bar shift=-3.0pt] table[x=x, y=cash]{data/clean/profiles/finh_600201.dat};
\addplot[fill=SoftRed!20, draw=SoftRed, pattern=north east lines, bar shift=3.0pt] table[x=x, y=ibd]{data/clean/profiles/finh_600201.dat};
\legend{货币资金（左轴）, 有息负债（左轴）}
\end{axis}
\begin{axis}[
  at={(finbot.south east)}, anchor=south east,
  width=13.75cm, height=3.10cm, scale only axis,
  axis y line*=right, axis x line*=none, xtick=\empty,
  ymin=-1.39, ymax=6.93,
  tick label style={font=\tiny},
  yticklabel style={font=\tiny}, scaled y ticks=false,
  legend style={font=\scriptsize, at={(0.5,-0.44)}, anchor=north,
                legend columns=1, draw=none, fill=none},
]
\addplot[OilBlue, mark=*, mark size=2.2pt, line width=1.3pt] table[x=i, y=ocf]{data/clean/profiles/fin_600201.dat};
\addplot[only marks, OilBlue, mark=*, mark size=3.2pt] table[x=x, y=ocf]{data/clean/profiles/finh_600201.dat};
\legend{经营活动现金流净额（右轴，亿元，蓝点）}
\end{axis}
\end{tikzpicture}

\captionof{figure}[生物股份（600201）近年主要财务指标]{生物股份（600201）近年主要财务指标（左上：营业收入与归母净利润；右上：ROE 与资产负债率；下：货币资金、有息负债为左轴柱状，经营活动现金流净额为其右轴的蓝点折线。
2021---2025 年为年报口径，2026H1 为 2026 年半年报/半年末，与其前一年报之间留出间隔、以斜纹或空心点区分；半年报数据未年化）}
\end{minipage}
\end{center}

\pfl{财务分析} 2026 年 6 月末资产负债率 16.5\%，较 2025 年末上升 0.4 个百分点（样本中位数 52.2\%，所处三级行业内按由低到高排第\zhnumber{4}位）。 2026 年 6 月末货币资金 16.09 亿元、短期债务（短期借款与一年内到期非流动负债）0.01 亿元、长期债务 1.00 亿元，有息负债合计 1.01 亿元，现金短债比 1,426.53 倍——覆盖充分。 净有息负债 -15.08 亿元，相当于其 2026 年上半年营业收入的 -199\%。 流动比率 4.21、速动比率 3.82，短期流动性无明显缺口。 2026 年上半年经营活动现金流净额 0.51 亿元，经营现金流与利润同向为正，内生资金可以支撑运营。 2026 年 6 月末在建工程 3.63 亿元，较上年末减少 0.02 亿元，在建投入在收缩。 2026 年 6 月末商誉 1.39 亿元，占归母净资产的 2.5\%，是净资产中最脆弱的一块。 综合看，现金短债比 1,426.53 倍，短期偿付压力可控；资产负债率低于样本中位数 35.7 个百分点；有息负债中长期债务占比更高，期限结构对短债滚动的压力较小。

\pfl{重大战略转型} 近三年公司的战略动作集中在：融资与资本运作（1 项）：2023 年 2 月，公司规划通过向特定对象发行募集资金投入动物 mRNA 疫苗及核酸药物开发项目（2.98 亿元）与动物 mRNA 疫苗及核酸药物生产车间建设项目（2.13 亿元）。 公告口径上，近三年（2023 年 9 月---2026 年 9 月）公司披露重大重组与并购类公告 0 条、对外投资与转型类公告 0 条、回购与股权激励类公告 40 条，合计公告 315 条。 第一大行业仍是“生物制药”，收入占比 92.7\%，与 2021 年年报相比没有方向性变化。 综合判断，报告期内公司未披露实质性的重组、并购或转型安排，资本开支以维持现有产能与业务为主。

\begin{center}\small\setlength{\tabcolsep}{4pt}
\captionof{table}[生物股份（600201）历史融资与资本运作]{生物股份（600201）历史融资与资本运作（据公司公告与公开报道整理；金额为披露值，含首次公开发行、再融资、债券、银行借款与重整投资等）}
\begin{adjustbox}{max width=\textwidth}
\begin{tabular}{@{}p{1.45cm}p{2.55cm}p{4.05cm}p{4.95cm}@{}}
\toprule
时间 & 方式 & 规模 & 用途与说明 \\
\midrule
1999-01 & 首次公开发行A股股票并在上海证券交易所上市（IPO） & 发行价 6.83 元/\allowbreak{}股；首发前总股本 6，\allowbreak{}200.25 万股，\allowbreak{}实际发行 3 & 招股公告日 1998 年 11 月 30 日 \\
2016-09 & 非公开发行A股股票（2015年度非公开发行） & 发行 4，\allowbreak{}032.26 万股，\allowbreak{}发行价格 31.00 元/\allowbreak{}股；募集资金总额 12.50 亿元 & 投资建设金宇生物科技产业园区项目一期工程。；定价基准日为第八届董事会第二十二次会议决议公告日（2016 年 2 月 25 日），\allowbreak{}底价 24.43 元/\allowbreak{}股；发行后实收资本（股本）变更为 6.13 亿元 \\
2021-07-07 & 非公开发行A股股票（2021年度非公开发行，后终止） & 拟非公开发行不超过 6，\allowbreak{}315.79 万股股票（含本数） & 2021 年 7 月 23 日 2021 年第一次临时股东大会审议通过 \\
2022-04-07 & 2021年度非公开发行终止 & 终止拟向内蒙古金宇生物控股有限公司非公开发行不超过 6,\allowbreak{}315.79 万股的计划 & 不适用。；终止原因为鉴于内外部环境变化等因素，\allowbreak{}综合公司整体规划 \\
2023-02-07 & 向特定对象发行A股股票（2023年度向特定对象发行） & 发行价格 8.29 元/\allowbreak{}股；拟募集资金总额不超过人民币 8.00 亿元（含本数）；发行对象为内蒙古金宇生物控股有限公司、\allowbreak{}张翀宇先生和张竞女士 & 拟投向动物 mRNA 疫苗及核酸药物开发项目 2.98 亿元、\allowbreak{}动物 mRNA 疫苗及核酸药物生产车间建设项目 2.13 亿元、\allowbreak{}补充流动资金 2.89 亿元。；2023 年 3 月 16 日 2023 年第二次临时股东大会审议通过；2023 年 5 月 9 日董事会审议通过调整方案；2024 年 1 月 12 日董事会审议通过调减募集资金总额暨调整方案 \\
2024-07-24 & 2023年度向特定对象发行取得中国证监会注册批复 & 中国证监会同意金宇生物批复自同意注册之日（2024 年 7 月 16 日）起 12 个月内有效 & 对应前次发行预案的 mRNA 疫苗及核酸药物开发、\allowbreak{}生产车间建设及补充流动资金项目。；公司于 2024 年 6 月 25 日公告该申请获得上海证券交易所审核通过 \\
2024-10-31 & 股份回购（2024年以集中竞价交易方式回购公司股份） & 拟回购股份数量占总股本比例 0.71\%\textasciitilde{}1.43\%，\allowbreak{}拟回购金额 8 & 用于股权激励及/\allowbreak{}或员工持股计划等（以公司回购方案公告为准）。；2026 年 5 月 19 日公司回购注销部分股票（暨减少注册资本）通知债权人 \\
2025-07-14 & 2023年度向特定对象发行终止 & 终止向特定对象发行 A 股股票 & 不适用 \\
2026-06 & 股份回购并取得银行股票回购专项贷款承诺函 & 回购价格不超过 20.74 元/\allowbreak{}股 & 用于股份回购；2026 年 6 月 13 日公司公告取得该回购专项贷款承诺函 \\
\bottomrule
\end{tabular}
\end{adjustbox}
\end{center}

\pfl{历史融投资情况} 从现金流量表看，2025 年公司偿还债务支付的现金 0.15 亿元、吸收权益性投资 0.07 亿元，筹资活动现金流净额 -3.07 亿元（2023 年为偿还债务 0.10 亿元）。 2026 年上半年筹资活动现金流净额为 -1.13 亿元（取得借款 1.00 亿元、偿还债务 0.20 亿元，未年化），已转为净流出，处于净还债状态。 2025 年分配股利、利润或偿付利息支付的现金合计 0.33 亿元。 公开披露的融资记录共 9 笔，工具类型以 IPO、定向增发、银行借款为主；金额与用途见下表。 其中规模最大的两笔为：2026 年 6 月，银行借款，回购价格不超过 20.74 元/股；2023 年 2 月，定向增发，发行价格 8.29 元/股，拟募集资金总额不超过人民币 8.00 亿元（含本数），发行对象为内蒙古金宇生物控股有限公司、张翀宇先生和张竞女士。 股本方面，近三年披露股本变动 4 次（其他 1 次）；总股本由 2021 年末的 11.26 亿股变为 11.20 亿股（-0.52\%）。 分红方面，近三年现金分红 4 次、合计每股派现 0.21 元，最近一次实施在 2025 年 12 月。 近三年“再融资”类公告 43 条，涉及定向增发、募集资金使用。 综合看，公司融资结构上以债务滚动为主、股权融资为补充；2025 年筹资活动现金流净额 -3.07 亿元，融资节奏仍以维持存量债务周转为主，与前述经营与投资现金流的方向基本一致。

\pfl{未来潜在融资需求分析} 2026 年 6 月末货币资金 16.09 亿元，而短期债务 0.01 亿元（现金短债比 1,426.53 倍），2026 年上半年经营现金流净额 0.51 亿元。账面货币资金可以覆盖短期债务，缺口不体现在静态偿付层面。 公司 2026 年 6 月末资产负债率 16.5\%、2026 年上半年 ROE 1.9\%（未年化），资产负债表尚有余量，融资需求不是“补血型”的刚性需求，而是配套在建项目投入的主动性安排。 公开披露的后续资金安排线索包括：股权融资线索：2026 年 4 月，公司未公告银行综合授信额度、公司债、中票或短融的发行计划，亦无新的定向增发预案；债务与授信安排：2026 年 6 月，贷款金额不超过 3.60 亿元（不超过回购金额的 90\%）。 资金需求还要叠加在建工程：期末在建工程 3.63 亿元、2026 年上半年购建固定资产等支出 0.44 亿元（未年化），这部分支出在周期底部通常会被主动放缓。 综合看，公司在静态偿付层面不存在缺口，外部融资更多是配合扩张节奏与并购整合的主动性安排，而非偿债压力驱动。




\subsubsection{中牧股份（600195）}
\label{co:600195}

\pfl{公司基本情况} 北京丰台的央企动保平台，全称中牧实业股份有限公司，1998 年 12 月由中国牧工商（集团）总公司独家发起以募集方式设立，1999 年 1 月 7 日在上海证券交易所上市；专业从事动物保健品和动物营养品研发、生产、销售及技术服务，设立 20 余个生产基地，产品覆盖兽用生物制品、兽用化药、饲料及饲料添加剂与饲料原料贸易，年营收超 60 亿元。2025 年归母净利润 1.71 亿元、同比增长 141.68\%，但扣非净利润仅 0.20 亿元、同比下降 67.42\%；2026 年上半年归母净利润 0.50 亿元、同比下降 39.04\%。 公司于 1999-01-07 在上交所首发上市，注册资本 10.21 亿元，总股本 10.21 亿股。 截至 2026 年 9 月 24 日收盘价 6.64 元，流通市值 67.8 亿元，近一年-6.0\%，流通市值在三级行业“动物保健Ⅲ”的 9 家公司中按由大到小排第\zhnumber{2}位，近一年涨跌幅高于全样本中位数（-10.1\%）4.1 个百分点。 股权结构方面，控股股东为中国牧工商集团有限公司；实际控制人为中国农业发展集团有限公司。

\pfl{历史股价} 自 2021 年 1 月至 2026 年 9 月，股价由 13.89 元变为 6.64 元，累计 -52.2\%，区间最高 15.92 元（2023 年 2 月）、最低 5.90 元（2026 年 6 月）；当前价相当于区间最高价的 41.7\%，为区间最低价的 1.13 倍。 月度收益的年化波动率 30.7\%，低于全样本中位数 37.4\%，在三级行业“动物保健Ⅲ”的 9 家公司中波动率由高到低排第\zhnumber{8}位。 把股价阶段与公司事件对齐看，2021 年 8 月 至 2023 年 2 月股价上涨+66.7\%，同期（2022 年 12 月）公司终止乾元浩分拆创业板上市后，公司未再公告新的资产分拆或重大资产重组方案。 从时间对应关系看，公司层面的资本运作与股权、风险事件解释了区间内多数急涨急跌，与行业指数的同步性相对有限。

\begin{center}
\begin{minipage}{\textwidth}\centering
\begin{tikzpicture}
\begin{axis}[
  width=12.6cm, height=3.9cm, scale only axis,
  xmin=2020.922, xmax=2026.888, ymin=5.49, ymax=17.03,
  xtick={2021,2022,2023,2024,2025,2026},
  yearticks,
  yticklabel style={font=\scriptsize},
  ylabel={元/股}, ylabel style={font=\scriptsize},
  grid=major, grid style={LightGray, line width=0.3pt},
  axis line style={InkGray, line width=0.5pt},
  every axis plot/.append style={line width=0.9pt},
]
\addplot[AgriGreen] table[x=x,y=y]{data/clean/profiles/px_600195.dat};
\addplot[SoftRed, dashed, line width=0.7pt] table[x=x,y=y]{data/clean/profiles/pxzz_600195.dat};
\addplot[only marks, mark=*, mark size=1.1pt, SoftRed] table[x=x,y=y]{data/clean/profiles/pxzz_600195.dat};
\end{axis}
\end{tikzpicture}

\captionof{figure}[中牧股份（600195）股价与周期骨架]{中牧股份（600195）月度收盘价（元/股）与 ZigZag 周期骨架（阈值 25\%，圆点为周期转折点）}
\end{minipage}
\end{center}

\pfl{过去周期分析} 以 25\% 的 ZigZag 阈值划分，样本区间内股价形成 2 个主升段与 3 个主跌段。 最大主跌段为 2023 年 2 月 至 2024 年 8 月，18 个月-61.5\%；最大主升段出现在 2021 年 8 月 至 2023 年 2 月，18 个月+66.7\%。 最近一次转折出现在 2026 年 6 月（下跌段自 2026 年 2 月起，4 个月-27.9\%），当前处于该段的延续之中。 公司股价与申万农林牧渔指数几乎同步见底，个股并未走出独立行情。 动保的需求由存栏量而非价格决定（第 \ref{sec:vet} 节）：猪价下行阶段龙头养殖企业的逆势扩产反而推升了疫苗与兽药需求，这正是动保股价与猪价“脱钩”的原因，公司 2026 年上半年实现盈利、半年 ROE 0.9\%（未年化），好于全样本中位数（0.75\%），下行周期对其报表的冲击小于同业。 综合区间形态看，该股单段涨跌幅度偏大、以趋势段为主（最大主升 +66.7\%、最大主跌 -61.5\%），年化波动率 30.7\% 与全样本中位数 37.4\% 相比波动小于样本中位数。

\pfl{盈利情况} 2026 年上半年营业收入 31.71 亿元（同比 +13.5\%），归母净利润 0.50 亿元，同比 -39.0\%。 以年报为趋势对照，2023---2025 年营业收入由 54.06 亿元变为 62.05 亿元（两年年均复合增速 +7.1\%），归母净利润由 4.03 亿元变为 1.71 亿元。 按半年报口径，2026 年上半年的 ROE 为 0.9\%（未年化）、销售净利率 2.3\%、毛利率 18.4\%，ROE 在“动物保健Ⅲ”9 家公司中按数值由高到低排第\zhnumber{8}位。 同期全样本半年 ROE 中位数 0.75\%，公司与之相差 0.13 个百分点（略好于中位数公司）。 公司披露的应对举措集中在：共设立 20 余个生产基地。 综合看，公司在报告期维持盈利（高于全样本半年 ROE 中位数 0.75\%），利润的可持续性取决于动物保健Ⅲ景气的延续与成本管控的执行。

\begin{center}\small\setlength{\tabcolsep}{3.4pt}
\captionof{table}[中牧股份（600195）关键财务指标]{中牧股份（600195）关键财务指标（合并报表口径；两个半年报期均未年化；现金短债比 $=$ 货币资金 $\div$ 短期债务）}
\begin{adjustbox}{max width=\textwidth}
\begin{tabular}{@{}lrrrrr@{}}
\toprule
指标 & 2023 年 & 2024 年 & 2025 年 & 2025 年半年报 & 2026 年半年报 \\
\midrule
营业收入（亿元） & 54.06 & 60.17 & 62.05 & 27.93 & 31.71 \\
归母净利润（亿元） & 4.03 & 0.71 & 1.71 & 0.82 & 0.50 \\
毛利率（\%） & 19.8 & 16.4 & 17.5 & 17.2 & 18.4 \\
ROE（\%） & 7.5 & 1.3 & 3.1 & 1.5 & 0.9 \\
资产负债率（\%） & 26.9 & 26.5 & 29.8 & 25.7 & 30.2 \\
经营活动现金流净额（亿元） & 3.72 & 4.20 & 11.58 & -2.32 & -2.42 \\
货币资金（亿元） & 14.93 & 10.91 & 18.73 & 8.96 & 16.65 \\
有息负债合计（亿元） & 13.64 & 10.18 & 10.46 & 11.17 & 11.12 \\
现金短债比（倍） & 3.23 & 3.09 & 2.51 & 1.94 & 1.97 \\
\bottomrule
\end{tabular}
\end{adjustbox}
\end{center}

\begin{center}
\begin{minipage}{\textwidth}\centering
\begin{tikzpicture}
\begin{groupplot}[
  group style={group size=2 by 1, horizontal sep=0.60cm},
  width=6.85cm, height=4.60cm,
  tick label style={font=\tiny},
  title style={font=\scriptsize\heiti},
  yticklabel style={font=\tiny},
  ylabel style={font=\tiny},
  grid=major, grid style={LightGray, line width=0.3pt},
  axis line style={InkGray, line width=0.5pt},
  enlarge x limits={abs=0.8},
  legend style={font=\scriptsize, at={(0.5,-0.20)}, anchor=north,
                legend columns=2, draw=none, fill=none, column sep=6pt},
]
\nextgroupplot[
  ybar, bar width=4pt,
  ymin=-6.20, ymax=69.49,
  xtick={0,1,2,3,4,5.7},
  xticklabels={2021,2022,2023,2024,2025,2026H1},
  ylabel={亿元},
]
\addplot[fill=AgriGreen!75, draw=AgriGreen, bar shift=-2.6pt] table[x=i, y=rev]{data/clean/profiles/fin_600195.dat};
\addplot[fill=SoftRed!60, draw=SoftRed, bar shift=2.6pt] table[x=i, y=ni]{data/clean/profiles/fin_600195.dat};
\addplot[fill=AgriGreen!30, draw=AgriGreen, pattern=north east lines, bar shift=-2.6pt] table[x=x, y=rev]{data/clean/profiles/finh_600195.dat};
\addplot[fill=SoftRed!20, draw=SoftRed, pattern=north east lines, bar shift=2.6pt] table[x=x, y=ni]{data/clean/profiles/finh_600195.dat};
\legend{营业收入（年/半年）, 归母净利润（年/半年）}
\nextgroupplot[
  ymin=-2.4, ymax=33.3,
  xtick={0,1,2,3,4,5.7},
  xticklabels={2021,2022,2023,2024,2025,2026H1},
  ylabel={\%},
]
\addplot[AgriGold, mark=*, mark size=1.3pt] table[x=i, y=roe]{data/clean/profiles/fin_600195.dat};
\addplot[OilBlue, mark=square*, mark size=1.3pt] table[x=i, y=debt]{data/clean/profiles/fin_600195.dat};
\addplot[only marks, AgriGold, mark=*, mark size=2.0pt] table[x=x, y=roe]{data/clean/profiles/finh_600195.dat};
\addplot[only marks, OilBlue, mark=square*, mark size=2.0pt] table[x=x, y=debt]{data/clean/profiles/finh_600195.dat};
\legend{ROE, 资产负债率}
\end{groupplot}
\end{tikzpicture}
\begin{tikzpicture}
\begin{axis}[
  name=finbot,
  width=13.75cm, height=3.10cm, scale only axis,
  ymin=-1.87, ymax=20.98,
  xtick={0,1,2,3,4,5.7},
  xticklabels={2021,2022,2023,2024,2025,2026H1},
  tick label style={font=\tiny},
  yticklabel style={font=\tiny},
  ylabel={亿元}, ylabel style={font=\tiny},
  ybar, bar width=4pt,
  grid=major, grid style={LightGray, line width=0.3pt},
  axis line style={InkGray, line width=0.5pt},
  enlarge x limits={abs=0.8},
  legend style={font=\scriptsize, at={(0.5,-0.26)}, anchor=north,
                legend columns=2, draw=none, fill=none, column sep=10pt},
]
\addplot[fill=AgriGreen!55, draw=AgriGreen, bar shift=-3.0pt] table[x=i, y=cash]{data/clean/profiles/fin_600195.dat};
\addplot[fill=SoftRed!45, draw=SoftRed, bar shift=3.0pt] table[x=i, y=ibd]{data/clean/profiles/fin_600195.dat};
\addplot[fill=AgriGreen!25, draw=AgriGreen, pattern=north east lines, bar shift=-3.0pt] table[x=x, y=cash]{data/clean/profiles/finh_600195.dat};
\addplot[fill=SoftRed!20, draw=SoftRed, pattern=north east lines, bar shift=3.0pt] table[x=x, y=ibd]{data/clean/profiles/finh_600195.dat};
\legend{货币资金（左轴）, 有息负债（左轴）}
\end{axis}
\begin{axis}[
  at={(finbot.south east)}, anchor=south east,
  width=13.75cm, height=3.10cm, scale only axis,
  axis y line*=right, axis x line*=none, xtick=\empty,
  ymin=-6.06, ymax=15.77,
  tick label style={font=\tiny},
  yticklabel style={font=\tiny}, scaled y ticks=false,
  legend style={font=\scriptsize, at={(0.5,-0.44)}, anchor=north,
                legend columns=1, draw=none, fill=none},
]
\addplot[OilBlue, mark=*, mark size=2.2pt, line width=1.3pt] table[x=i, y=ocf]{data/clean/profiles/fin_600195.dat};
\addplot[only marks, OilBlue, mark=*, mark size=3.2pt] table[x=x, y=ocf]{data/clean/profiles/finh_600195.dat};
\legend{经营活动现金流净额（右轴，亿元，蓝点）}
\end{axis}
\end{tikzpicture}

\captionof{figure}[中牧股份（600195）近年主要财务指标]{中牧股份（600195）近年主要财务指标（左上：营业收入与归母净利润；右上：ROE 与资产负债率；下：货币资金、有息负债为左轴柱状，经营活动现金流净额为其右轴的蓝点折线。
2021---2025 年为年报口径，2026H1 为 2026 年半年报/半年末，与其前一年报之间留出间隔、以斜纹或空心点区分；半年报数据未年化）}
\end{minipage}
\end{center}

\pfl{财务分析} 2026 年 6 月末资产负债率 30.2\%，较 2025 年末上升 0.4 个百分点（样本中位数 52.2\%，所处三级行业内按由低到高排第\zhnumber{5}位）。 2026 年 6 月末货币资金 16.65 亿元、短期债务（短期借款与一年内到期非流动负债）8.44 亿元、长期债务 2.69 亿元，有息负债合计 11.12 亿元，现金短债比 1.97 倍——覆盖充分。 净有息负债 -5.52 亿元，相当于其 2026 年上半年营业收入的 -17\%。 流动比率 1.77、速动比率 1.35，短期流动性无明显缺口。 2026 年上半年经营活动现金流净额 -2.42 亿元，净利润为正但经营现金流为负，利润的现金含量偏低。 2026 年 6 月末在建工程 0.58 亿元，较上年末增加 0.03 亿元，仍有在建投入。 综合看，现金短债比 1.97 倍，短期偿付压力可控；资产负债率低于样本中位数 22.0 个百分点；有息负债中短期债务占比更高，期限结构仍以滚动续作为主。 公司披露的诉讼与资产冻结风险为：2026 年 6 月，第二大股东中国华农资产经营有限公司所持公司 1，668.35 万股（占总股本 1.63\%）处于冻结状态。

\pfl{重大战略转型} 近三年公司的战略动作集中在：并购与重组（1 项）：2022 年 12 月，公司未再公告新的资产分拆或重大资产重组方案。 公告口径上，近三年（2023 年 9 月---2026 年 9 月）公司披露重大重组与并购类公告 2 条、对外投资与转型类公告 0 条、回购与股权激励类公告 0 条，合计公告 261 条。 综合判断，报告期内公司有 2 条与战略相关的公告落地（重大重组与并购 2 条、对外投资与转型 0 条），资本运作与主业调整之间存在直接关联。

\begin{center}\small\setlength{\tabcolsep}{4pt}
\captionof{table}[中牧股份（600195）历史融资与资本运作]{中牧股份（600195）历史融资与资本运作（据公司公告与公开报道整理；金额为披露值，含首次公开发行、再融资、债券、银行借款与重整投资等）}
\begin{adjustbox}{max width=\textwidth}
\begin{tabular}{@{}p{1.45cm}p{2.55cm}p{4.05cm}p{4.95cm}@{}}
\toprule
时间 & 方式 & 规模 & 用途与说明 \\
\midrule
1999-01 & 首次公开发行A股股票并在上海证券交易所上市（IPO） & 发行价 6.28 元/\allowbreak{}股；首发前总股本 18，\allowbreak{}000.00 万股，\allowbreak{}实际发行 8 & 限责任公司，\allowbreak{}招股公告日 1998 年 11 月 16 日 \\
2013-10 & 非公开发行A股股票（定增） & 发行 3，\allowbreak{}980 万股 & 公司在认购邀请书规定的有效申报时间内共收到 6 家投资者发出的有效申购报价单；新增股份为有限售条件流通股 \\
2015-09-18 & 重大资产重组（停牌筹划并终止） & --- & 不适用。；公司因筹划股权收购事项股票自 2015 年 9 月 18 日起进入重大资产重组停牌程序 \\
2016-02 & 公开发行公司债券（2016年第一期） & 过 12.00 亿元的公司债券；本期基础发行规模 8.00 亿元 & 用于拓宽融资渠道、\allowbreak{}优化债务结构及补充营运资金（以募集说明书为准）。路 66 号 4 号楼；2016 年 2 月 24 日披露发行公告 \\
2019-12-10 & 公开发行公司债券（预案） & 公司董事会审议通过公司债券发行预案，\allowbreak{}拟申请公开发行公司债券 & 为进一步拓宽融资渠道、\allowbreak{}优化债务结构、\allowbreak{}满足资金需求。；该预案公告编号为临 2019-055 \\
2021-07-23 & 分拆控股子公司乾元浩生物股份有限公司至创业板上市 & --- & 不适用。；公司启动分拆乾元浩至创业板上市事项 \\
2022-12-30 & 终止分拆控股子公司乾元浩至创业板上市 & 终止分拆控股子公司乾元浩生物股份有限公司至创业板上市 & 不适用。；公司于 2022 年 12 月 30 日召开第八届董事会 2022 年第二十次临时会议、\allowbreak{}第八届监事会 2022 年第四次临时会议审议通过终止筹划控股子公司分拆上市编号临 2022-045 \\
\bottomrule
\end{tabular}
\end{adjustbox}
\end{center}

\pfl{历史融投资情况} 从现金流量表看，2025 年公司取得借款收到的现金 8.01 亿元、偿还债务支付的现金 7.40 亿元、吸收权益性投资 0.72 亿元，筹资活动现金流净额 0.36 亿元（2023 年为取得借款 13.46 亿元、偿还债务 3.61 亿元）。 2026 年上半年筹资活动现金流净额为 0.55 亿元（取得借款 5.54 亿元、偿还债务 4.87 亿元，未年化），仍处于净融入状态。 2025 年分配股利、利润或偿付利息支付的现金合计 0.93 亿元。 公开披露的融资记录共 7 笔，工具类型以 IPO、定向增发、公司债为主；金额与用途见下表。 其中规模最大的两笔为：2016 年 2 月，公司债，过 12.00 亿元的公司债券，本期基础发行规模 8.00 亿元；2013 年 10 月，定向增发，发行 3，980 万股。 股本方面，近三年披露股本变动 3 次；总股本由 2021 年末的 10.16 亿股变为 10.21 亿股（+0.55\%）。 分红方面，近三年现金分红 5 次、合计每股派现 0.52 元，最近一次实施在 2026 年 7 月。 近三年“再融资”类公告 3 条，涉及可转债。 综合看，公司融资结构上股权与债务融资并举；2025 年筹资活动现金流净额 0.36 亿元，年度内仍为净融入，与前述经营与投资现金流的方向基本一致。

\pfl{未来潜在融资需求分析} 2026 年 6 月末货币资金 16.65 亿元，而短期债务 8.44 亿元（现金短债比 1.97 倍），2026 年上半年经营现金流净额 -2.42 亿元。账面货币资金可以覆盖短期债务，缺口不体现在静态偿付层面。 公司 2026 年 6 月末资产负债率 30.2\%、2026 年上半年 ROE 0.9\%（未年化），资产端与负债端都没有明显压力，融资需求以相机抉择的扩张型用途为主。 公开披露的后续资金安排线索包括：股权融资线索：2026 年 8 月，公司未公告新的定向增发预案、配股、可转债或公司债发行计划，亦未公告 2026 年度银行综合授信额度；债务与授信安排：2019 年 12 月，公司曾于 2019 年 12 月披露公司债券发行预案。 资金需求还要叠加在建工程：期末在建工程 0.58 亿元、2026 年上半年购建固定资产等支出 0.77 亿元（未年化），这部分支出在周期底部通常会被主动放缓。 静态偿付不存在缺口，后续融资更可能服务于产能或并购安排，而不是填补流动性窟窿。




\subsubsection{科前生物（688526）}
\label{co:688526}

\pfl{公司基本情况} 科前生物成立于 2016 年，武汉的兽用生物制品企业，全称武汉科前生物股份有限公司，2020 年 9 月 22 日在上海证券交易所科创板上市；专注于兽用生物制品研发、生产、销售及动物防疫技术服务，主要产品是非国家强制免疫猪用疫苗和禽用疫苗，可生产 85 个兽用生物制品类产品，是我国兽用生物制品行业内产品种类最丰富的企业之一，并已向禽用疫苗、宠物疫苗、诊断试剂等领域拓展。公司无控股股东，实际控制人为何启盖、吴斌、方六荣、吴美洲四人（一致行动人协议）。2025 年归母净利润 4.19 亿元、同比增长 9.51\%；2026 年上半年归母净利润 1.16 亿元、同比下降 47.37\%；2023 年一次定增撤回、一次定增被上交所终止审核。 公司于 2020-09-22 在上交所科创板首发上市，注册资本 4.66 亿元，总股本 4.66 亿股。 截至 2026 年 9 月 24 日收盘价 12.18 元，流通市值 56.8 亿元，近一年-31.0\%，流通市值在三级行业“动物保健Ⅲ”的 9 家公司中按由大到小排第\zhnumber{3}位，近一年涨跌幅低于全样本中位数（-10.1\%）20.9 个百分点。 按 2025 年年报披露的行业构成，收入占比前三位为兽用生物制品 93.8\%；其他 4.8\%；其他(补充) 1.5\%。 与 2021 年年报相比，第一大行业口径由“兽用生物制造行业”（98.2\%）变为“兽用生物制品”（93.8\%）；两期披露的分类口径有调整，占比差异中同时包含分类因素。

\pfl{历史股价} 以 2021 年 1 月为起点、2026 年 9 月为终点，股价由 38.08 元变为 12.18 元，累计 -68.0\%，区间最高 41.46 元（2021 年 4 月）、最低 12.10 元（2026 年 6 月）；当前价相当于区间最高价的 29.4\%，已接近区间最低价（1.01 倍）。 月度收益的年化波动率 37.0\%，低于全样本中位数 37.4\%，在三级行业“动物保健Ⅲ”的 9 家公司中波动率由高到低排第\zhnumber{6}位。 把股价阶段与公司事件对齐看，2023 年 2 月 至 2024 年 8 月股价下跌-60.4\%，同期（2024 年 5 月）本次投资不构成关联交易，亦不构成重大资产重组。 把这些阶段与公司事件对齐后可以看到，几轮大涨大跌多由公司自身的资本运作与风险事件触发，行业景气只解释了其中一部分。

\begin{center}
\begin{minipage}{\textwidth}\centering
\begin{tikzpicture}
\begin{axis}[
  width=12.6cm, height=3.9cm, scale only axis,
  xmin=2020.922, xmax=2026.888, ymin=11.25, ymax=44.36,
  xtick={2021,2022,2023,2024,2025,2026},
  yearticks,
  yticklabel style={font=\scriptsize},
  ylabel={元/股}, ylabel style={font=\scriptsize},
  grid=major, grid style={LightGray, line width=0.3pt},
  axis line style={InkGray, line width=0.5pt},
  every axis plot/.append style={line width=0.9pt},
]
\addplot[AgriGreen] table[x=x,y=y]{data/clean/profiles/px_688526.dat};
\addplot[SoftRed, dashed, line width=0.7pt] table[x=x,y=y]{data/clean/profiles/pxzz_688526.dat};
\addplot[only marks, mark=*, mark size=1.1pt, SoftRed] table[x=x,y=y]{data/clean/profiles/pxzz_688526.dat};
\end{axis}
\end{tikzpicture}

\captionof{figure}[科前生物（688526）股价与周期骨架]{科前生物（688526）月度收盘价（元/股）与 ZigZag 周期骨架（阈值 25\%，圆点为周期转折点）}
\end{minipage}
\end{center}

\pfl{过去周期分析} 以 25\% 的 ZigZag 阈值划分，样本区间内股价形成 2 个主升段与 3 个主跌段。 最大主跌段为 2023 年 2 月 至 2024 年 8 月，18 个月-60.4\%；最大主升段出现在 2022 年 4 月 至 2023 年 2 月，10 个月+73.1\%。 最近一次转折出现在 2026 年 6 月（下跌段自 2025 年 8 月起，10 个月-34.2\%），当前处于该段的延续之中。 公司股价与申万农林牧渔指数几乎同步见底，个股并未走出独立行情。 动保的需求由存栏量而非价格决定（第 \ref{sec:vet} 节）：猪价下行阶段龙头养殖企业的逆势扩产反而推升了疫苗与兽药需求，这正是动保股价与猪价“脱钩”的原因，公司 2026 年上半年实现盈利、半年 ROE 2.8\%（未年化），好于全样本中位数（0.75\%），下行周期对其报表的冲击小于同业。 综合区间形态看，该股单段涨跌幅度偏大、以趋势段为主（最大主升 +73.1\%、最大主跌 -60.4\%），年化波动率 37.0\% 与全样本中位数 37.4\% 相比波动与样本中位数接近。

\pfl{盈利情况} 2026 年上半年营业收入 3.74 亿元（同比 -23.3\%），归母净利润 1.16 亿元，同比 -47.4\%。 以年报为趋势对照，2023---2025 年营业收入由 10.64 亿元变为 9.50 亿元（两年年均复合增速 -5.5\%），归母净利润由 3.96 亿元变为 4.19 亿元。 按半年报口径，2026 年上半年的 ROE 为 2.8\%（未年化）、销售净利率 30.7\%、毛利率 59.4\%，ROE 在“动物保健Ⅲ”9 家公司中按数值由高到低排第\zhnumber{6}位。 同期全样本半年 ROE 中位数 0.75\%，公司与之相差 2.02 个百分点（略好于中位数公司）。 从公司披露的原因看，业绩变动主要来自：2026 年 4 月，公司商誉账面余额 330.75 万元，系收购武汉诸乐田源生态农业有限公司形成。 针对上述压力，公司披露的应对举措包括：2024 年 5 月，公司终止首发募投项目“动物生物制品车间技改项目”，将结余募集资金 2.87 亿元转投“高级别动物疫苗产业化基地建设项目（一期）”。 综合看，公司在报告期维持盈利（高于全样本半年 ROE 中位数 0.75\%），利润的可持续性取决于动物保健Ⅲ景气的延续与成本管控的执行。

\begin{center}\small\setlength{\tabcolsep}{3.4pt}
\captionof{table}[科前生物（688526）关键财务指标]{科前生物（688526）关键财务指标（合并报表口径；两个半年报期均未年化；现金短债比 $=$ 货币资金 $\div$ 短期债务）}
\begin{adjustbox}{max width=\textwidth}
\begin{tabular}{@{}lrrrrr@{}}
\toprule
指标 & 2023 年 & 2024 年 & 2025 年 & 2025 年半年报 & 2026 年半年报 \\
\midrule
营业收入（亿元） & 10.64 & 9.42 & 9.50 & 4.87 & 3.74 \\
归母净利润（亿元） & 3.96 & 3.82 & 4.19 & 2.20 & 1.16 \\
毛利率（\%） & 72.1 & 63.8 & 67.0 & 67.3 & 59.4 \\
ROE（\%） & 10.8 & 9.9 & 10.3 & 5.4 & 2.8 \\
资产负债率（\%） & 17.8 & 12.3 & 11.5 & 11.4 & 10.7 \\
经营活动现金流净额（亿元） & 5.38 & 3.93 & 4.69 & 1.74 & 1.16 \\
货币资金（亿元） & 5.82 & 5.43 & 1.77 & 5.95 & 1.70 \\
有息负债合计（亿元） & 2.94 & 0.07 & 0.07 & 0.07 & 0.07 \\
现金短债比（倍） & 2.02 & 402.35 & 110.03 & 426.69 & 96.61 \\
\bottomrule
\end{tabular}
\end{adjustbox}
\end{center}

\begin{center}
\begin{minipage}{\textwidth}\centering
\begin{tikzpicture}
\begin{groupplot}[
  group style={group size=2 by 1, horizontal sep=0.60cm},
  width=6.85cm, height=4.60cm,
  tick label style={font=\tiny},
  title style={font=\scriptsize\heiti},
  yticklabel style={font=\tiny},
  ylabel style={font=\tiny},
  grid=major, grid style={LightGray, line width=0.3pt},
  axis line style={InkGray, line width=0.5pt},
  enlarge x limits={abs=0.8},
  legend style={font=\scriptsize, at={(0.5,-0.20)}, anchor=north,
                legend columns=2, draw=none, fill=none, column sep=6pt},
]
\nextgroupplot[
  ybar, bar width=4pt,
  ymin=-1.10, ymax=12.35,
  xtick={0,1,2,3,4,5.7},
  xticklabels={2021,2022,2023,2024,2025,2026H1},
  ylabel={亿元},
]
\addplot[fill=AgriGreen!75, draw=AgriGreen, bar shift=-2.6pt] table[x=i, y=rev]{data/clean/profiles/fin_688526.dat};
\addplot[fill=SoftRed!60, draw=SoftRed, bar shift=2.6pt] table[x=i, y=ni]{data/clean/profiles/fin_688526.dat};
\addplot[fill=AgriGreen!30, draw=AgriGreen, pattern=north east lines, bar shift=-2.6pt] table[x=x, y=rev]{data/clean/profiles/finh_688526.dat};
\addplot[fill=SoftRed!20, draw=SoftRed, pattern=north east lines, bar shift=2.6pt] table[x=x, y=ni]{data/clean/profiles/finh_688526.dat};
\legend{营业收入（年/半年）, 归母净利润（年/半年）}
\nextgroupplot[
  ymin=-1.6, ymax=22.3,
  xtick={0,1,2,3,4,5.7},
  xticklabels={2021,2022,2023,2024,2025,2026H1},
  ylabel={\%},
]
\addplot[AgriGold, mark=*, mark size=1.3pt] table[x=i, y=roe]{data/clean/profiles/fin_688526.dat};
\addplot[OilBlue, mark=square*, mark size=1.3pt] table[x=i, y=debt]{data/clean/profiles/fin_688526.dat};
\addplot[only marks, AgriGold, mark=*, mark size=2.0pt] table[x=x, y=roe]{data/clean/profiles/finh_688526.dat};
\addplot[only marks, OilBlue, mark=square*, mark size=2.0pt] table[x=x, y=debt]{data/clean/profiles/finh_688526.dat};
\legend{ROE, 资产负债率}
\end{groupplot}
\end{tikzpicture}
\begin{tikzpicture}
\begin{axis}[
  name=finbot,
  width=13.75cm, height=3.10cm, scale only axis,
  ymin=-0.58, ymax=6.52,
  xtick={0,1,2,3,4,5.7},
  xticklabels={2021,2022,2023,2024,2025,2026H1},
  tick label style={font=\tiny},
  yticklabel style={font=\tiny},
  ylabel={亿元}, ylabel style={font=\tiny},
  ybar, bar width=4pt,
  grid=major, grid style={LightGray, line width=0.3pt},
  axis line style={InkGray, line width=0.5pt},
  enlarge x limits={abs=0.8},
  legend style={font=\scriptsize, at={(0.5,-0.26)}, anchor=north,
                legend columns=2, draw=none, fill=none, column sep=10pt},
]
\addplot[fill=AgriGreen!55, draw=AgriGreen, bar shift=-3.0pt] table[x=i, y=cash]{data/clean/profiles/fin_688526.dat};
\addplot[fill=SoftRed!45, draw=SoftRed, bar shift=3.0pt] table[x=i, y=ibd]{data/clean/profiles/fin_688526.dat};
\addplot[fill=AgriGreen!25, draw=AgriGreen, pattern=north east lines, bar shift=-3.0pt] table[x=x, y=cash]{data/clean/profiles/finh_688526.dat};
\addplot[fill=SoftRed!20, draw=SoftRed, pattern=north east lines, bar shift=3.0pt] table[x=x, y=ibd]{data/clean/profiles/finh_688526.dat};
\legend{货币资金（左轴）, 有息负债（左轴）}
\end{axis}
\begin{axis}[
  at={(finbot.south east)}, anchor=south east,
  width=13.75cm, height=3.10cm, scale only axis,
  axis y line*=right, axis x line*=none, xtick=\empty,
  ymin=-1.40, ymax=6.99,
  tick label style={font=\tiny},
  yticklabel style={font=\tiny}, scaled y ticks=false,
  legend style={font=\scriptsize, at={(0.5,-0.44)}, anchor=north,
                legend columns=1, draw=none, fill=none},
]
\addplot[OilBlue, mark=*, mark size=2.2pt, line width=1.3pt] table[x=i, y=ocf]{data/clean/profiles/fin_688526.dat};
\addplot[only marks, OilBlue, mark=*, mark size=3.2pt] table[x=x, y=ocf]{data/clean/profiles/finh_688526.dat};
\legend{经营活动现金流净额（右轴，亿元，蓝点）}
\end{axis}
\end{tikzpicture}

\captionof{figure}[科前生物（688526）近年主要财务指标]{科前生物（688526）近年主要财务指标（左上：营业收入与归母净利润；右上：ROE 与资产负债率；下：货币资金、有息负债为左轴柱状，经营活动现金流净额为其右轴的蓝点折线。
2021---2025 年为年报口径，2026H1 为 2026 年半年报/半年末，与其前一年报之间留出间隔、以斜纹或空心点区分；半年报数据未年化）}
\end{minipage}
\end{center}

\pfl{财务分析} 2026 年 6 月末资产负债率 10.7\%，较 2025 年末下降 0.8 个百分点（样本中位数 52.2\%，所处三级行业内按由低到高排第\zhnumber{2}位）。 2026 年 6 月末货币资金 1.70 亿元、短期债务（短期借款与一年内到期非流动负债）0.02 亿元、长期债务 0.05 亿元，有息负债合计 0.07 亿元，现金短债比 96.61 倍——覆盖充分。 净有息负债 -1.63 亿元，相当于其 2026 年上半年营业收入的 -44\%。 流动比率 7.82、速动比率 7.44，短期流动性无明显缺口。 2026 年上半年经营活动现金流净额 1.16 亿元，经营现金流与利润同向为正，内生资金可以支撑运营。 综合看，现金短债比 96.61 倍，短期偿付压力可控；资产负债率低于样本中位数 41.5 个百分点；有息负债中长期债务占比更高，期限结构对短债滚动的压力较小。

\pfl{重大战略转型} 从公开披露看，公司的战略推进集中在：其他动作（2 项）：2026 年上半年，较上年同期增加 6.56 个百分点，在营收下滑背景下继续加大研发投入；主业转型（1 项）：2023 年 9 月，公司实际控制人由陈焕春、金梅林、何启盖、吴斌、方六荣、吴美洲、叶长发七人变更为何启盖、吴斌、方六荣、吴美洲四人本次实际控制人变更不会对公司生产经营产生重大影响。 公告口径上，近三年（2023 年 9 月---2026 年 9 月）公司披露重大重组与并购类公告 0 条、对外投资与转型类公告 1 条、回购与股权激励类公告 35 条，合计公告 343 条。 主营结构的口径变化已经落到报表上：2021 年年报的第一大行业为“兽用生物制造行业”（98.2\%），2025 年年报为“兽用生物制品”（93.8\%）；公司在这两期的分部划分方式有过调整。 综合判断，报告期内公司有 1 条与战略相关的公告落地（重大重组与并购 0 条、对外投资与转型 1 条），资本运作与主业调整之间存在直接关联。

\begin{center}\small\setlength{\tabcolsep}{4pt}
\captionof{table}[科前生物（688526）历史融资与资本运作]{科前生物（688526）历史融资与资本运作（据公司公告与公开报道整理；金额为披露值，含首次公开发行、再融资、债券、银行借款与重整投资等）}
\begin{adjustbox}{max width=\textwidth}
\begin{tabular}{@{}p{1.45cm}p{2.55cm}p{4.05cm}p{4.95cm}@{}}
\toprule
时间 & 方式 & 规模 & 用途与说明 \\
\midrule
2020-09 & 首次公开发行A股股票并在上海证券交易所科创板上市（IPO） & 发行价 11.69 元/\allowbreak{}股；首发前总股本 36，\allowbreak{}000.00 万股，\allowbreak{}实际发行 10 & 投向动物生物制品产业化建设项目（项目投资总额 9.03 亿元、\allowbreak{}募集资金承诺投资额 5.78 亿元）、\allowbreak{}动物生物制品车间技改项目（募集资金承诺投资额 2.87 亿元）、\allowbreak{}科研创新项目及补充流动资金；公司无超募资金。；募集资金于 2020 年 9 月 17 日到位 \\
2022-10-28 & 向特定对象发行A股股票（2022年度向特定对象发行） & 拟向特定对象发行股票数量不超过 876.60 万股（含本数） & 用于“高级别动物疫苗产业化基地建设项目（动物生物安全实验室）”。；2022 年 10 月 28 日董事会、\allowbreak{}2022 年 11 月 14 日股东大会审议通过；2023 年 2 月 27 日第三届董事会第二十四次会议、\allowbreak{}第三届监事会第二十二次会议审议通过 2022 年度向特定对象发行 A 股股票预案（修订稿） \\
2023-06-26 & 终止2022年度向特定对象发行A股股票并撤回申请文件 & 终止拟发行不超过 876.60 万股、\allowbreak{}募集资金总额不超过 1.62 亿元的发行计划 & 不适用。；公司于 2023 年 6 月 26 日召开第三届董事会第二十八次会议和第三届监事会第二十六次会议 \\
2024-05-17 & 变更部分募集资金投资项目 & 拟将首发募投项目“动物生物制品车间技改项目”的募集资金 2.87 亿元（截至 2024 年 4 月 30 日）变更用于新项目“高级别动物疫苗产业化基地建设项目（一期）” & 新项目预计 2028 年 10 月竣工（最终以实际开展情况为准）。；2024 年 5 月 17 日第三届董事会第四十次会议、\allowbreak{}第三届监事会第三十二次会议审议通过 \\
2025-12-31 & 首次公开发行募集资金使用进展 & 截至 2025 年 12 月 31 日，\allowbreak{}首发募集资金净额 11.42 亿元 & 按募投项目计划投入 \\
\bottomrule
\end{tabular}
\end{adjustbox}
\end{center}

\pfl{历史融投资情况} 从现金流量表看，2025 年公司，筹资活动现金流净额 -1.84 亿元（2023 年为取得借款 2.37 亿元、偿还债务 2.22 亿元）。 2026 年上半年筹资活动现金流净额为 -3.00 亿元（取得借款 — 亿元、偿还债务 — 亿元，未年化），已转为净流出，处于净还债状态。 2025 年分配股利、利润或偿付利息支付的现金合计 1.99 亿元。 公开披露的融资记录共 5 笔，工具类型以 IPO、定向增发为主；金额与用途见下表。 其中规模最大的两笔为：2024 年 5 月，变更部分募集资金投资，拟将首发募投项目“动物生物制品车间技改项目”的募集资金 2.87 亿元（截至 2024 年 4 月 30 日）变更用于新项目“高级别动物疫苗产业化基地建设项目（一期）”；2020 年 9 月，IPO，发行价 11.69 元/股，首发前总股本 36，000.00 万股，实际发行 10。 股本方面，近三年披露股本变动 7 次（股份回购 2 次、限售股份上市 1 次、其他 1 次）；总股本由 2021 年末的 4.65 亿股变为 4.66 亿股（+0.21\%）。 分红方面，近三年现金分红 4 次、合计每股派现 1.69 元，最近一次实施在 2025 年 12 月。 综合看，公司融资结构上股权与债务融资并举；2025 年筹资活动现金流净额 -1.84 亿元，融资节奏仍以维持存量债务周转为主，与前述经营与投资现金流的方向基本一致。

\pfl{未来潜在融资需求分析} 2026 年 6 月末货币资金 1.70 亿元，而短期债务 0.02 亿元（现金短债比 96.61 倍），2026 年上半年经营现金流净额 1.16 亿元。账面货币资金可以覆盖短期债务，缺口不体现在静态偿付层面。 公司 2026 年 6 月末资产负债率 10.7\%、2026 年上半年 ROE 2.8\%（未年化），账面杠杆不高、盈利能力一般，短期内没有必须依靠外部融资才能维持的经营缺口；后续融资更可能指向技改、扩产或产业并购。 公开披露的后续资金安排线索包括：股权融资线索：2026 年 8 月，公司未公告新的定向增发预案、可转债、公司债或银行综合授信额度安排；资本开支安排：2025 年 12 月，首发募集资金累计投入进度为 84.5\%。 综合看，公司在静态偿付层面不存在缺口，外部融资更多是配合扩张节奏与并购整合的主动性安排，而非偿债压力驱动。




\subsubsection{瑞普生物（300119）}
\label{co:300119}

\pfl{公司基本情况} 天津自贸试验区的动物保健企业，全称瑞普生物股份有限公司，2008 年 5 月 19 日由天津瑞普生物技术集团有限公司整体变更设立，2010 年 9 月在深圳证券交易所创业板上市；所属行业为医药制造类，产品体系涵盖兽用生物制品、兽用药物制剂、功能性添加剂与兽用原料药，并参股瑞派宠物医院（持股 9.09\%）布局宠物医疗。控股股东、实际控制人为自然人李守军（持股 36.04\%，任董事长）。2025 年归母净利润 4.01 亿元、同比增长 33.18\%；2026 年上半年归母净利润 2.04 亿元、同比下降 20.43\%；2026 年 3 月董事会提请股东会授权以简易程序向特定对象发行融资总额不超过 3 亿元的股票。 公司于 2010-09-17 在深交所创业板首发上市，注册资本 4.65 亿元，总股本 4.66 亿股。 截至 2026 年 9 月 24 日收盘价 16.66 元，流通市值 56.4 亿元，近一年-20.2\%，流通市值在三级行业“动物保健Ⅲ”的 9 家公司中按由大到小排第\zhnumber{4}位，近一年涨跌幅低于全样本中位数（-10.1\%）10.1 个百分点。 按 2026 年半年报披露的行业构成，收入占比前两位为动保业务 75.0\%；宠物供应链 25.0\%。 与 2021 年年报相比，第一大行业口径由“兽用生物制品”（53.3\%）变为“动保业务”（75.0\%）；两期披露的分类口径有调整，占比差异中同时包含分类因素。

\pfl{历史股价} 本报告样本区间（2021 年 1 月 至 2026 年 9 月）内，股价由 24.52 元变为 16.66 元，累计 -32.1\%，区间最高 38.05 元（2021 年 5 月）、最低 12.09 元（2024 年 8 月）；当前价相当于区间最高价的 43.8\%，为区间最低价的 1.38 倍。 月度收益的年化波动率 40.1\%，高于全样本中位数 37.4\%，在三级行业“动物保健Ⅲ”的 9 家公司中波动率由高到低排第\zhnumber{3}位。

\begin{center}
\begin{minipage}{\textwidth}\centering
\begin{tikzpicture}
\begin{axis}[
  width=12.6cm, height=3.9cm, scale only axis,
  xmin=2020.922, xmax=2026.888, ymin=11.24, ymax=40.71,
  xtick={2021,2022,2023,2024,2025,2026},
  yearticks,
  yticklabel style={font=\scriptsize},
  ylabel={元/股}, ylabel style={font=\scriptsize},
  grid=major, grid style={LightGray, line width=0.3pt},
  axis line style={InkGray, line width=0.5pt},
  every axis plot/.append style={line width=0.9pt},
]
\addplot[AgriGreen] table[x=x,y=y]{data/clean/profiles/px_300119.dat};
\addplot[SoftRed, dashed, line width=0.7pt] table[x=x,y=y]{data/clean/profiles/pxzz_300119.dat};
\addplot[only marks, mark=*, mark size=1.1pt, SoftRed] table[x=x,y=y]{data/clean/profiles/pxzz_300119.dat};
\end{axis}
\end{tikzpicture}

\captionof{figure}[瑞普生物（300119）股价与周期骨架]{瑞普生物（300119）月度收盘价（元/股）与 ZigZag 周期骨架（阈值 25\%，圆点为周期转折点）}
\end{minipage}
\end{center}

\pfl{过去周期分析} 以 25\% 的 ZigZag 阈值划分，样本区间内股价形成 3 个主升段与 3 个主跌段。 最大主跌段为 2021 年 5 月 至 2022 年 4 月，11 个月-59.9\%；最大主升段出现在 2024 年 8 月 至 2025 年 5 月，9 个月+103.3\%。 最近一次转折出现在 2026 年 6 月（下跌段自 2025 年 5 月起，13 个月-41.3\%），当前处于该段的延续之中。 公司股价较申万农林牧渔指数提前约 21 个月见底，是板块中较早定价周期反转的公司之一。 动保的需求由存栏量而非价格决定（第 \ref{sec:vet} 节）：猪价下行阶段龙头养殖企业的逆势扩产反而推升了疫苗与兽药需求，这正是动保股价与猪价“脱钩”的原因，公司 2026 年上半年实现盈利、半年 ROE 4.3\%（未年化），好于全样本中位数（0.75\%），下行周期对其报表的冲击小于同业。 综合区间形态看，该股单段涨跌幅度偏大、以趋势段为主（最大主升 +103.3\%、最大主跌 -59.9\%），年化波动率 40.1\% 与全样本中位数 37.4\% 相比波动与样本中位数接近。

\pfl{盈利情况} 2026 年上半年营业收入 17.20 亿元（同比 +0.8\%），归母净利润 2.04 亿元，同比 -20.4\%。 以年报为趋势对照，2023---2025 年营业收入由 27.09 亿元变为 33.98 亿元（两年年均复合增速 +12.0\%），归母净利润由 4.54 亿元变为 4.01 亿元。 按半年报口径，2026 年上半年的 ROE 为 4.3\%（未年化）、销售净利率 12.9\%、毛利率 39.7\%，ROE 在“动物保健Ⅲ”9 家公司中按数值由高到低排第\zhnumber{4}位。 按同一口径，三级行业“动物保健Ⅲ”9 家公司的半年 ROE 中位数为 2.95\%，该公司高于其中位数 1.32 个百分点。 综合看，公司在报告期维持盈利（高于全样本半年 ROE 中位数 0.75\%），利润的可持续性取决于动物保健Ⅲ景气的延续与成本管控的执行。

\begin{center}\small\setlength{\tabcolsep}{3.4pt}
\captionof{table}[瑞普生物（300119）关键财务指标]{瑞普生物（300119）关键财务指标（合并报表口径；两个半年报期均未年化；现金短债比 $=$ 货币资金 $\div$ 短期债务）}
\begin{adjustbox}{max width=\textwidth}
\begin{tabular}{@{}lrrrrr@{}}
\toprule
指标 & 2023 年 & 2024 年 & 2025 年 & 2025 年半年报 & 2026 年半年报 \\
\midrule
营业收入（亿元） & 27.09 & 30.70 & 33.98 & 17.08 & 17.20 \\
归母净利润（亿元） & 4.54 & 3.01 & 4.01 & 2.57 & 2.04 \\
毛利率（\%） & 49.7 & 41.4 & 41.3 & 42.1 & 39.7 \\
ROE（\%） & 10.2 & 6.5 & 8.8 & 5.6 & 4.3 \\
资产负债率（\%） & 27.2 & 37.1 & 32.3 & 36.0 & 34.4 \\
经营活动现金流净额（亿元） & 4.10 & 7.36 & 2.98 & 1.35 & 1.94 \\
货币资金（亿元） & 3.39 & 4.87 & 6.51 & 5.33 & 5.32 \\
有息负债合计（亿元） & 8.97 & 13.81 & 14.60 & 16.22 & 16.16 \\
现金短债比（倍） & 0.59 & 0.57 & 0.60 & 0.56 & 0.42 \\
\bottomrule
\end{tabular}
\end{adjustbox}
\end{center}

\begin{center}
\begin{minipage}{\textwidth}\centering
\begin{tikzpicture}
\begin{groupplot}[
  group style={group size=2 by 1, horizontal sep=0.60cm},
  width=6.85cm, height=4.60cm,
  tick label style={font=\tiny},
  title style={font=\scriptsize\heiti},
  yticklabel style={font=\tiny},
  ylabel style={font=\tiny},
  grid=major, grid style={LightGray, line width=0.3pt},
  axis line style={InkGray, line width=0.5pt},
  enlarge x limits={abs=0.8},
  legend style={font=\scriptsize, at={(0.5,-0.20)}, anchor=north,
                legend columns=2, draw=none, fill=none, column sep=6pt},
]
\nextgroupplot[
  ybar, bar width=4pt,
  ymin=-3.40, ymax=38.06,
  xtick={0,1,2,3,4,5.7},
  xticklabels={2021,2022,2023,2024,2025,2026H1},
  ylabel={亿元},
]
\addplot[fill=AgriGreen!75, draw=AgriGreen, bar shift=-2.6pt] table[x=i, y=rev]{data/clean/profiles/fin_300119.dat};
\addplot[fill=SoftRed!60, draw=SoftRed, bar shift=2.6pt] table[x=i, y=ni]{data/clean/profiles/fin_300119.dat};
\addplot[fill=AgriGreen!30, draw=AgriGreen, pattern=north east lines, bar shift=-2.6pt] table[x=x, y=rev]{data/clean/profiles/finh_300119.dat};
\addplot[fill=SoftRed!20, draw=SoftRed, pattern=north east lines, bar shift=2.6pt] table[x=x, y=ni]{data/clean/profiles/finh_300119.dat};
\legend{营业收入（年/半年）, 归母净利润（年/半年）}
\nextgroupplot[
  ymin=-3.0, ymax=40.9,
  xtick={0,1,2,3,4,5.7},
  xticklabels={2021,2022,2023,2024,2025,2026H1},
  ylabel={\%},
]
\addplot[AgriGold, mark=*, mark size=1.3pt] table[x=i, y=roe]{data/clean/profiles/fin_300119.dat};
\addplot[OilBlue, mark=square*, mark size=1.3pt] table[x=i, y=debt]{data/clean/profiles/fin_300119.dat};
\addplot[only marks, AgriGold, mark=*, mark size=2.0pt] table[x=x, y=roe]{data/clean/profiles/finh_300119.dat};
\addplot[only marks, OilBlue, mark=square*, mark size=2.0pt] table[x=x, y=debt]{data/clean/profiles/finh_300119.dat};
\legend{ROE, 资产负债率}
\end{groupplot}
\end{tikzpicture}
\begin{tikzpicture}
\begin{axis}[
  name=finbot,
  width=13.75cm, height=3.10cm, scale only axis,
  ymin=-1.62, ymax=18.10,
  xtick={0,1,2,3,4,5.7},
  xticklabels={2021,2022,2023,2024,2025,2026H1},
  tick label style={font=\tiny},
  yticklabel style={font=\tiny},
  ylabel={亿元}, ylabel style={font=\tiny},
  ybar, bar width=4pt,
  grid=major, grid style={LightGray, line width=0.3pt},
  axis line style={InkGray, line width=0.5pt},
  enlarge x limits={abs=0.8},
  legend style={font=\scriptsize, at={(0.5,-0.26)}, anchor=north,
                legend columns=2, draw=none, fill=none, column sep=10pt},
]
\addplot[fill=AgriGreen!55, draw=AgriGreen, bar shift=-3.0pt] table[x=i, y=cash]{data/clean/profiles/fin_300119.dat};
\addplot[fill=SoftRed!45, draw=SoftRed, bar shift=3.0pt] table[x=i, y=ibd]{data/clean/profiles/fin_300119.dat};
\addplot[fill=AgriGreen!25, draw=AgriGreen, pattern=north east lines, bar shift=-3.0pt] table[x=x, y=cash]{data/clean/profiles/finh_300119.dat};
\addplot[fill=SoftRed!20, draw=SoftRed, pattern=north east lines, bar shift=3.0pt] table[x=x, y=ibd]{data/clean/profiles/finh_300119.dat};
\legend{货币资金（左轴）, 有息负债（左轴）}
\end{axis}
\begin{axis}[
  at={(finbot.south east)}, anchor=south east,
  width=13.75cm, height=3.10cm, scale only axis,
  axis y line*=right, axis x line*=none, xtick=\empty,
  ymin=-1.91, ymax=9.56,
  tick label style={font=\tiny},
  yticklabel style={font=\tiny}, scaled y ticks=false,
  legend style={font=\scriptsize, at={(0.5,-0.44)}, anchor=north,
                legend columns=1, draw=none, fill=none},
]
\addplot[OilBlue, mark=*, mark size=2.2pt, line width=1.3pt] table[x=i, y=ocf]{data/clean/profiles/fin_300119.dat};
\addplot[only marks, OilBlue, mark=*, mark size=3.2pt] table[x=x, y=ocf]{data/clean/profiles/finh_300119.dat};
\legend{经营活动现金流净额（右轴，亿元，蓝点）}
\end{axis}
\end{tikzpicture}

\captionof{figure}[瑞普生物（300119）近年主要财务指标]{瑞普生物（300119）近年主要财务指标（左上：营业收入与归母净利润；右上：ROE 与资产负债率；下：货币资金、有息负债为左轴柱状，经营活动现金流净额为其右轴的蓝点折线。
2021---2025 年为年报口径，2026H1 为 2026 年半年报/半年末，与其前一年报之间留出间隔、以斜纹或空心点区分；半年报数据未年化）}
\end{minipage}
\end{center}

\pfl{财务分析} 2026 年 6 月末资产负债率 34.4\%，较 2025 年末上升 2.1 个百分点（样本中位数 52.2\%，所处三级行业内按由低到高排第\zhnumber{7}位）。 2026 年 6 月末货币资金 5.32 亿元、短期债务（短期借款与一年内到期非流动负债）12.68 亿元、长期债务 3.49 亿元，有息负债合计 16.16 亿元，现金短债比 0.42 倍——覆盖偏紧。 净有息负债 10.85 亿元，相当于其 2026 年上半年营业收入的 63\%。 流动比率 1.62、速动比率 1.38，短期偿债能力处于正常区间。 2026 年上半年经营活动现金流净额 1.94 亿元，经营现金流与利润同向为正，内生资金可以支撑运营。 2026 年 6 月末在建工程 0.81 亿元，较上年末增加 0.49 亿元，仍有在建投入。 2026 年 6 月末商誉 4.46 亿元，占归母净资产的 8.7\%，是净资产中最脆弱的一块。 综合看，账面货币资金对有息负债与短期债务的覆盖不足，滚动偿付对新增融资的依赖度较高；资产负债率低于样本中位数 17.8 个百分点；有息负债中短期债务占比更高，期限结构仍以滚动续作为主。 公司披露的股东与质押风险为：2026 年 6 月，控股股东、实际控制人李守军所持公司股份中质押 4，720.00 万股。

\pfl{重大战略转型} 近三年公司的战略动作集中在：其他动作（3 项）：2026 年 3 月，公司通过参股瑞派宠物医院管理股份有限公司（9.09\%）并推进瑞普生物、瑞派宠物医院、中瑞供应链“三瑞齐发”战略合作；2025 年 9 月，公司与福建圣维生物科技有限公司合资设立内蒙古瑞圣生物科技有限公司；并购与重组（1 项）：2024 年 4 月，公司以 2.30 亿元自有或自筹资金收购保定市收骏科技有限公司 100\% 股权（关联交易）。 公告口径上，近三年（2023 年 9 月---2026 年 9 月）公司披露重大重组与并购类公告 13 条、对外投资与转型类公告 3 条、回购与股权激励类公告 34 条，合计公告 388 条。 主营结构的口径变化已经落到报表上：2021 年年报的第一大行业为“兽用生物制品”（53.3\%），2026 年半年报为“动保业务”（75.0\%）；公司在这两期的分部划分方式有过调整。 综合判断，报告期内公司有 16 条与战略相关的公告落地（重大重组与并购 13 条、对外投资与转型 3 条），资本运作与主业调整之间存在直接关联。

\begin{center}\small\setlength{\tabcolsep}{4pt}
\captionof{table}[瑞普生物（300119）历史融资与资本运作]{瑞普生物（300119）历史融资与资本运作（据公司公告与公开报道整理；金额为披露值，含首次公开发行、再融资、债券、银行借款与重整投资等）}
\begin{adjustbox}{max width=\textwidth}
\begin{tabular}{@{}p{1.45cm}p{2.55cm}p{4.05cm}p{4.95cm}@{}}
\toprule
时间 & 方式 & 规模 & 用途与说明 \\
\midrule
2010-09 & 首次公开发行A股股票并在深圳证券交易所创业板上市（IPO） & 发行价 60.00 元/\allowbreak{}股；首发前总股本 5，\allowbreak{}554.80 万股，\allowbreak{}实际发行 1 & 限责任公司，\allowbreak{}发行方式为网下询价配售与网上定价发行 \\
2016-08 & 创业板非公开发行A股股票（2015年度非公开发行） & 发行 1，\allowbreak{}534.72 万股，\allowbreak{}发行价 14.92 元/\allowbreak{}股；募集资金总额 2.29 亿元 & 市莞城区可园南路一号；发行对象包括李守军、\allowbreak{}梁武、\allowbreak{}李旭东、\allowbreak{}鲍恩东等（以发行情况报告书为准） \\
2021-11 & 向特定对象发行股票（2021年定增） & 发行 6，\allowbreak{}398.47 万股，\allowbreak{}发行价 20.88 元/\allowbreak{}股；募集资金总额 13.36 亿元 & 投向公司向特定对象发行股票募集资金投资项目；其中部分募集资金已于 2025 年变更用途。；新增股份于 2021 年 11 月 19 日在深圳证券交易所上市 \\
2024-04 & 收购保定市收骏科技有限公司100\%股权（关联交易，自有资金或自筹资金） & 交易金额 2.30 亿元 & 2024 年 3 月 28 日第五届董事会第十三次会议、\allowbreak{}第五届监事会第十次会议审议通过；2024 年 4 月公司公告完成工商变更登记 \\
2025-04-01 & 向控股子公司提供财务资助 & 向控股子公司中瑞供应链提供额度不超过人民币 5，\allowbreak{}000.00 万元的借款 & 支持控股子公司经营周转。；2025 年 4 月 1 日第五届董事会第二十二次会议及第五届监事会第十八次会议审议通过编号 2025-024 \\
2025-06-13 & 向控股子公司提供财务资助 & 向中瑞供应链提供额度不超过人民币 1.20 亿元的借款，\allowbreak{}额度范围内可循环使用 & 支持控股子公司经营周转。；2025 年 6 月 13 日第五届董事会第二十四次（临时）会议审议通过编号 2025-047 \\
2025-09-28 & 变更部分募集资金用途 & 变更 2021 年向特定对象发行股票部分募集资金投资项目用途 & 新项目为“生物制造产业化工程建设项目”。；2025 年 9 月 28 日第五届董事会第二十七次（临时）会议审议通过编号 2025-076 \\
2026-03-27 & 提请股东会授权董事会以简易程序向特定对象发行股票 & 融资总额不超过人民币 3.00 亿元且不超过最近一年末净资产 20\% 的中国境内上市人民币普通股（A 股） & 募集资金具体用途由董事会在授权范围内确定 \\
\bottomrule
\end{tabular}
\end{adjustbox}
\end{center}

\pfl{历史融投资情况} 从现金流量表看，2025 年公司取得借款收到的现金 11.86 亿元、偿还债务支付的现金 10.59 亿元、吸收权益性投资 0.01 亿元，筹资活动现金流净额 -2.50 亿元（2023 年为取得借款 7.28 亿元、偿还债务 7.31 亿元）。 2026 年上半年筹资活动现金流净额为 -0.74 亿元（取得借款 8.25 亿元、偿还债务 6.70 亿元，未年化），已转为净流出，处于净还债状态。 2025 年分配股利、利润或偿付利息支付的现金合计 2.25 亿元。 公开披露的融资记录共 8 笔，工具类型以 IPO、定向增发为主；金额与用途见下表。 其中规模最大的两笔为：2010 年 9 月，IPO，发行价 60.00 元/股，首发前总股本 5，554.80 万股，实际发行 1；2026 年 3 月，定向增发，融资总额不超过人民币 3.00 亿元且不超过最近一年末净资产 20\% 的中国境内上市人民币普通股（A 股）。 股本方面，近三年披露股本变动 8 次（股份回购 2 次、其他 2 次、激励股份解禁 1 次）；总股本由 2021 年末的 4.68 亿股变为 4.66 亿股（-0.53\%）。 分红方面，近三年现金分红 4 次、合计每股派现 1.40 元，最近一次实施在 2025 年 12 月。 近三年“再融资”类公告 4 条，涉及定向增发。 综合看，公司融资结构上以债务滚动为主、股权融资为补充；2025 年筹资活动现金流净额 -2.50 亿元，融资节奏仍以维持存量债务周转为主，与前述经营与投资现金流的方向基本一致。

\pfl{未来潜在融资需求分析} 2026 年 6 月末货币资金 5.32 亿元，而短期债务 12.68 亿元（现金短债比 0.42 倍），2026 年上半年经营现金流净额 1.94 亿元。按“短期债务减去账面货币资金”的滚动偿付口径，静态缺口约 7.36 亿元；上半年盈利或接近盈亏平衡，该缺口不随经营亏损进一步扩大。 按第 \ref{sec:financing-need} 节的三档划分，公司属于第二档“保壳型”：近三年重大重组与并购类公告 13 条，2026 年上半年归母净利润 2.04 亿元，融资需求更多体现为“资产置换 + 维持上市地位”，而不是扩产。 公开披露的后续资金安排线索包括：股权融资线索：2025 年 9 月，公司变更 2021 年定增部分募集资金用途用于“生物制造产业化工程建设项目”；2026 年 3 月，公司董事会提请股东会授权以简易程序向特定对象发行股票。 资金需求还要叠加在建工程：期末在建工程 0.81 亿元、2026 年上半年购建固定资产等支出 0.97 亿元（未年化），这部分支出在周期底部通常会被主动放缓。 综合看，公司的外部资金需求以营运资金周转与在建项目投入为主，滚动偿付口径下的静态缺口约 7.36 亿元，能否落地取决于授信续作与发行窗口的配合。




\subsubsection{回盛生物（300871）}
\label{co:300871}

\pfl{公司基本情况} 武汉的兽药企业，全称武汉回盛生物科技股份有限公司，2002 年 1 月 25 日成立，2020 年 8 月 24 日在深圳证券交易所创业板上市；主营兽用药品（化药制剂、原料药、中药制剂）及饲料与添加剂，是农业产业化国家重点龙头企业，截至 2025 年 12 月 31 日已取得兽药批准文号 317 个，产品应用于生猪、家禽、水产、反刍及宠物等领域。控股股东为武汉统盛投资有限公司（持股 41.03\%）。2025 年归母净利润 2.53 亿元、同比增长 1356.01\%（2024 年为亏损 0.20 亿元）；2026 年上半年归母净利润 2.06 亿元、同比增长 75.59\%。2026 年 4 月推出逾两年的 2024 年度定增方案到期失效，5 月被深交所终止审核。 公司于 2020-08-24 在深交所创业板首发上市，注册资本 2.02 亿元，总股本 1.66 亿股。 截至 2026 年 9 月 24 日收盘价 20.89 元，流通市值 42.2 亿元，近一年+1.2\%，流通市值在三级行业“动物保健Ⅲ”的 9 家公司中按由大到小排第\zhnumber{5}位，近一年涨跌幅高于全样本中位数（-10.1\%）11.2 个百分点。 按 2025 年年报披露的行业构成，收入占比前两位为兽用原料药及制剂 95.5\%；其他 4.5\%。

\pfl{历史股价} 以 2021 年 1 月为起点、2026 年 9 月为终点，股价由 44.00 元变为 20.89 元，累计 -52.5\%，区间最高 52.92 元（2021 年 4 月）、最低 8.47 元（2024 年 8 月）；当前价相当于区间最高价的 39.5\%，为区间最低价的 2.47 倍。 月度收益的年化波动率 54.4\%，高于全样本中位数 37.4\%，在三级行业“动物保健Ⅲ”的 9 家公司中波动率由高到低排第\zhnumber{2}位。

\begin{center}
\begin{minipage}{\textwidth}\centering
\begin{tikzpicture}
\begin{axis}[
  width=12.6cm, height=3.9cm, scale only axis,
  xmin=2020.922, xmax=2026.888, ymin=7.88, ymax=56.62,
  xtick={2021,2022,2023,2024,2025,2026},
  yearticks,
  yticklabel style={font=\scriptsize},
  ylabel={元/股}, ylabel style={font=\scriptsize},
  grid=major, grid style={LightGray, line width=0.3pt},
  axis line style={InkGray, line width=0.5pt},
  every axis plot/.append style={line width=0.9pt},
]
\addplot[AgriGreen] table[x=x,y=y]{data/clean/profiles/px_300871.dat};
\addplot[SoftRed, dashed, line width=0.7pt] table[x=x,y=y]{data/clean/profiles/pxzz_300871.dat};
\addplot[only marks, mark=*, mark size=1.1pt, SoftRed] table[x=x,y=y]{data/clean/profiles/pxzz_300871.dat};
\end{axis}
\end{tikzpicture}

\captionof{figure}[回盛生物（300871）股价与周期骨架]{回盛生物（300871）月度收盘价（元/股）与 ZigZag 周期骨架（阈值 25\%，圆点为周期转折点）}
\end{minipage}
\end{center}

\pfl{过去周期分析} 以 25\% 的 ZigZag 阈值划分，样本区间内股价形成 4 个主升段与 4 个主跌段。 最大主跌段为 2021 年 4 月 至 2022 年 4 月，12 个月-69.4\%；最大主升段出现在 2024 年 8 月 至 2026 年 1 月，17 个月+247.8\%。 最近一次转折出现在 2026 年 7 月（下跌段自 2026 年 1 月起，6 个月-29.7\%），当前处于该段的延续之中。 公司股价较申万农林牧渔指数提前约 21 个月见底，是板块中较早定价周期反转的公司之一。 动保的需求由存栏量而非价格决定（第 \ref{sec:vet} 节）：猪价下行阶段龙头养殖企业的逆势扩产反而推升了疫苗与兽药需求，这正是动保股价与猪价“脱钩”的原因，公司 2026 年上半年实现盈利、半年 ROE 9.1\%（未年化），好于全样本中位数（0.75\%），下行周期对其报表的冲击小于同业。 综合区间形态看，该股单段涨跌幅度偏大、以趋势段为主（最大主升 +247.8\%、最大主跌 -69.4\%），年化波动率 54.4\% 与全样本中位数 37.4\% 相比波动明显大于样本中位数。

\pfl{盈利情况} 2026 年上半年营业收入 8.65 亿元（同比 +5.2\%），归母净利润 2.06 亿元，同比 +75.6\%。 以年报为趋势对照，2023---2025 年营业收入由 10.20 亿元变为 16.81 亿元（两年年均复合增速 +28.4\%），归母净利润由 0.17 亿元变为 2.53 亿元。 按半年报口径，2026 年上半年的 ROE 为 9.1\%（未年化）、销售净利率 23.9\%、毛利率 35.7\%，ROE 在“动物保健Ⅲ”9 家公司中按数值由高到低排第\zhnumber{2}位。 按同一口径，三级行业“动物保健Ⅲ”9 家公司的半年 ROE 中位数为 2.95\%，该公司高于其中位数 6.11 个百分点。 综合看，公司在报告期维持盈利（高于全样本半年 ROE 中位数 0.75\%），利润的可持续性取决于动物保健Ⅲ景气的延续与成本管控的执行。

\begin{center}\small\setlength{\tabcolsep}{3.4pt}
\captionof{table}[回盛生物（300871）关键财务指标]{回盛生物（300871）关键财务指标（合并报表口径；两个半年报期均未年化；现金短债比 $=$ 货币资金 $\div$ 短期债务）}
\begin{adjustbox}{max width=\textwidth}
\begin{tabular}{@{}lrrrrr@{}}
\toprule
指标 & 2023 年 & 2024 年 & 2025 年 & 2025 年半年报 & 2026 年半年报 \\
\midrule
营业收入（亿元） & 10.20 & 12.00 & 16.81 & 8.22 & 8.65 \\
归母净利润（亿元） & 0.17 & -0.20 & 2.53 & 1.17 & 2.06 \\
毛利率（\%） & 20.9 & 16.3 & 31.0 & 27.6 & 35.7 \\
ROE（\%） & 1.1 & -1.4 & 13.4 & 7.0 & 9.1 \\
资产负债率（\%） & 49.3 & 51.4 & 27.3 & 30.4 & 31.8 \\
经营活动现金流净额（亿元） & 0.93 & 0.44 & 2.38 & 1.24 & 1.28 \\
货币资金（亿元） & 2.94 & 3.26 & 3.10 & 2.36 & 3.59 \\
有息负债合计（亿元） & 8.42 & 8.89 & 2.29 & 3.50 & 4.64 \\
现金短债比（倍） & 1.22 & 0.98 & 1.52 & 0.75 & 1.77 \\
\bottomrule
\end{tabular}
\end{adjustbox}
\end{center}

\begin{center}
\begin{minipage}{\textwidth}\centering
\begin{tikzpicture}
\begin{groupplot}[
  group style={group size=2 by 1, horizontal sep=0.60cm},
  width=6.85cm, height=4.60cm,
  tick label style={font=\tiny},
  title style={font=\scriptsize\heiti},
  yticklabel style={font=\tiny},
  ylabel style={font=\tiny},
  grid=major, grid style={LightGray, line width=0.3pt},
  axis line style={InkGray, line width=0.5pt},
  enlarge x limits={abs=0.8},
  legend style={font=\scriptsize, at={(0.5,-0.20)}, anchor=north,
                legend columns=2, draw=none, fill=none, column sep=6pt},
]
\nextgroupplot[
  ybar, bar width=4pt,
  ymin=-1.90, ymax=18.85,
  xtick={0,1,2,3,4,5.7},
  xticklabels={2021,2022,2023,2024,2025,2026H1},
  ylabel={亿元},
]
\addplot[fill=AgriGreen!75, draw=AgriGreen, bar shift=-2.6pt] table[x=i, y=rev]{data/clean/profiles/fin_300871.dat};
\addplot[fill=SoftRed!60, draw=SoftRed, bar shift=2.6pt] table[x=i, y=ni]{data/clean/profiles/fin_300871.dat};
\addplot[fill=AgriGreen!30, draw=AgriGreen, pattern=north east lines, bar shift=-2.6pt] table[x=x, y=rev]{data/clean/profiles/finh_300871.dat};
\addplot[fill=SoftRed!20, draw=SoftRed, pattern=north east lines, bar shift=2.6pt] table[x=x, y=ni]{data/clean/profiles/finh_300871.dat};
\legend{营业收入（年/半年）, 归母净利润（年/半年）}
\nextgroupplot[
  ymin=-5.6, ymax=56.7,
  xtick={0,1,2,3,4,5.7},
  xticklabels={2021,2022,2023,2024,2025,2026H1},
  ylabel={\%},
]
\addplot[AgriGold, mark=*, mark size=1.3pt] table[x=i, y=roe]{data/clean/profiles/fin_300871.dat};
\addplot[OilBlue, mark=square*, mark size=1.3pt] table[x=i, y=debt]{data/clean/profiles/fin_300871.dat};
\addplot[only marks, AgriGold, mark=*, mark size=2.0pt] table[x=x, y=roe]{data/clean/profiles/finh_300871.dat};
\addplot[only marks, OilBlue, mark=square*, mark size=2.0pt] table[x=x, y=debt]{data/clean/profiles/finh_300871.dat};
\legend{ROE, 资产负债率}
\end{groupplot}
\end{tikzpicture}
\begin{tikzpicture}
\begin{axis}[
  name=finbot,
  width=13.75cm, height=3.10cm, scale only axis,
  ymin=-1.07, ymax=11.96,
  xtick={0,1,2,3,4,5.7},
  xticklabels={2021,2022,2023,2024,2025,2026H1},
  tick label style={font=\tiny},
  yticklabel style={font=\tiny},
  ylabel={亿元}, ylabel style={font=\tiny},
  ybar, bar width=4pt,
  grid=major, grid style={LightGray, line width=0.3pt},
  axis line style={InkGray, line width=0.5pt},
  enlarge x limits={abs=0.8},
  legend style={font=\scriptsize, at={(0.5,-0.26)}, anchor=north,
                legend columns=2, draw=none, fill=none, column sep=10pt},
]
\addplot[fill=AgriGreen!55, draw=AgriGreen, bar shift=-3.0pt] table[x=i, y=cash]{data/clean/profiles/fin_300871.dat};
\addplot[fill=SoftRed!45, draw=SoftRed, bar shift=3.0pt] table[x=i, y=ibd]{data/clean/profiles/fin_300871.dat};
\addplot[fill=AgriGreen!25, draw=AgriGreen, pattern=north east lines, bar shift=-3.0pt] table[x=x, y=cash]{data/clean/profiles/finh_300871.dat};
\addplot[fill=SoftRed!20, draw=SoftRed, pattern=north east lines, bar shift=3.0pt] table[x=x, y=ibd]{data/clean/profiles/finh_300871.dat};
\legend{货币资金（左轴）, 有息负债（左轴）}
\end{axis}
\begin{axis}[
  at={(finbot.south east)}, anchor=south east,
  width=13.75cm, height=3.10cm, scale only axis,
  axis y line*=right, axis x line*=none, xtick=\empty,
  ymin=-0.62, ymax=3.10,
  tick label style={font=\tiny},
  yticklabel style={font=\tiny}, scaled y ticks=false,
  legend style={font=\scriptsize, at={(0.5,-0.44)}, anchor=north,
                legend columns=1, draw=none, fill=none},
]
\addplot[OilBlue, mark=*, mark size=2.2pt, line width=1.3pt] table[x=i, y=ocf]{data/clean/profiles/fin_300871.dat};
\addplot[only marks, OilBlue, mark=*, mark size=3.2pt] table[x=x, y=ocf]{data/clean/profiles/finh_300871.dat};
\legend{经营活动现金流净额（右轴，亿元，蓝点）}
\end{axis}
\end{tikzpicture}

\captionof{figure}[回盛生物（300871）近年主要财务指标]{回盛生物（300871）近年主要财务指标（左上：营业收入与归母净利润；右上：ROE 与资产负债率；下：货币资金、有息负债为左轴柱状，经营活动现金流净额为其右轴的蓝点折线。
2021---2025 年为年报口径，2026H1 为 2026 年半年报/半年末，与其前一年报之间留出间隔、以斜纹或空心点区分；半年报数据未年化）}
\end{minipage}
\end{center}

\pfl{财务分析} 2026 年 6 月末资产负债率 31.8\%，较 2025 年末上升 4.4 个百分点（样本中位数 52.2\%，所处三级行业内按由低到高排第\zhnumber{6}位）。 2026 年 6 月末货币资金 3.59 亿元、短期债务（短期借款与一年内到期非流动负债）2.02 亿元、长期债务 2.62 亿元，有息负债合计 4.64 亿元，现金短债比 1.77 倍——覆盖充分。 净有息负债 1.05 亿元，相当于其 2026 年上半年营业收入的 12\%。 流动比率 2.05、速动比率 1.67，短期流动性无明显缺口。 2026 年上半年经营活动现金流净额 1.28 亿元，经营现金流与利润同向为正，内生资金可以支撑运营。 2026 年 6 月末在建工程 1.58 亿元，较上年末增加 1.00 亿元，仍有在建投入。 综合看，现金短债比 1.77 倍，短期偿付压力可控；资产负债率低于样本中位数 20.4 个百分点；有息负债中长期债务占比更高，期限结构对短债滚动的压力较小。 公司披露的股东与质押风险为：2026 年 6 月，控股股东武汉统盛投资有限公司所持公司股份中质押 420.00 万股；2024 年 4 月，可能涉及股票质押作为并购贷的担保措施。

\pfl{重大战略转型} 从公开披露看，公司的战略推进集中在：融资与资本运作（2 项）：2026 年 4 月，2021 年可转债募投项目中包含“宠物制剂综合生产线建设项目”（拟投入 9，000.00 万元）；2024 年 4 月，公司拟向楚盛投资及实际控制人发行股票募集不超过 2.50 亿元全部补充流动资金；股权与股东回报（1 项）：2026 年 6 月，公司实际控制人、董事长张卫元提议回购公司股份。 公告口径上，近三年（2023 年 9 月---2026 年 9 月）公司披露重大重组与并购类公告 2 条、对外投资与转型类公告 3 条、回购与股权激励类公告 36 条，合计公告 497 条。 第一大行业仍是“兽用原料药及制剂”，收入占比 95.5\%，与 2021 年年报相比没有方向性变化。 综合判断，报告期内公司有 5 条与战略相关的公告落地（重大重组与并购 2 条、对外投资与转型 3 条），资本运作与主业调整之间存在直接关联。

\begin{center}\small\setlength{\tabcolsep}{4pt}
\captionof{table}[回盛生物（300871）历史融资与资本运作]{回盛生物（300871）历史融资与资本运作（据公司公告与公开报道整理；金额为披露值，含首次公开发行、再融资、债券、银行借款与重整投资等）}
\begin{adjustbox}{max width=\textwidth}
\begin{tabular}{@{}p{1.45cm}p{2.55cm}p{4.05cm}p{4.95cm}@{}}
\toprule
时间 & 方式 & 规模 & 用途与说明 \\
\midrule
2020-08 & 首次公开发行A股股票并在深圳证券交易所创业板上市（IPO） & 发行价 33.61 元/\allowbreak{}股；首发前总股本 8，\allowbreak{}280.70 万股，\allowbreak{}实际发行 2 & 实际募集资金超出发行人预计募集资金。份有限公司 \\
2021-12-17 & 向不特定对象发行可转换公司债券（回盛转债，债券代码123132） & 发行 700.00 万张，\allowbreak{}每张面值 100.00 元 & 年产 1，\allowbreak{}000 吨泰乐菌素和年产 600 吨泰万菌素生产线扩建项目（拟投入 2.85 亿元 \\
2024-04-03 & 向特定对象发行A股股票（2024年度向特定对象发行，后终止） & 拟募集资金总额不超过人民币 2.50 亿元（含本数） & 全部用于补充流动资金。；2024 年 4 月 3 日第三届董事会第十次会议、\allowbreak{}第三届监事会第八次会议审议通过 \\
2024-04-25 & 向金融机构申请综合授信额度（关联担保） & 公司及子公司 2024 年度拟向银行等金融机构申请授信额度不超过人民币 10.00 亿元（含本数） & 融资方式包括银行流动资金贷款、\allowbreak{}固定资产投资贷款、\allowbreak{}保函、\allowbreak{}信用证、\allowbreak{}银行承兑汇票、\allowbreak{}各类商业票据开立及贴现、\allowbreak{}保理等授信业务。；公司公告明确该授信额度不等同于公司实际融资金额 \\
2026-04-17 & 2024年度向特定对象发行A股股票方案到期失效 & 终止拟募集不超过 2.50 亿元（补充流动资金）的发行计划 & 不适用。；公司 2026 年第二次临时股东会审议未通过延长公司 2024 年度向特定对象发行 A 股股票股东会决议有效期和提请股东会延长授权董事会及其授权人士全权办理公司向特定对象发行 A 股股票相关事宜有效期 \\
2026-05 & 深圳证券交易所终止审核2024年度向特定对象发行 & 终止审核公司向特定对象发行股票申请 & 不适用。；公司于 2026 年 4 月 17 日召开的 2026 年第二次临时股东会审议未通过延长决议有效期及授权的议案后 \\
2026-06-05 & 实际控制人、董事长提议回购公司股份 & 公司实际控制人、\allowbreak{}董事长张卫元先生于 2026 年 6 月 5 日出具提议回购公司股份 & 基于对公司价值的判断和未来发展前景的信心提议回购。；该事项处于提议阶段 \\
\bottomrule
\end{tabular}
\end{adjustbox}
\end{center}

\pfl{历史融投资情况} 从现金流量表看，2025 年公司取得借款收到的现金 2.20 亿元、偿还债务支付的现金 3.68 亿元、吸收权益性投资 0.00 亿元，筹资活动现金流净额 -1.60 亿元（2023 年为取得借款 2.11 亿元、偿还债务 2.44 亿元）。 2026 年上半年筹资活动现金流净额为 2.07 亿元（取得借款 4.24 亿元、偿还债务 2.10 亿元，未年化），仍处于净融入状态。 2025 年分配股利、利润或偿付利息支付的现金合计 0.50 亿元。 公开披露的融资记录共 7 笔，工具类型以 IPO、可转债、定向增发、银行授信为主；金额与用途见下表。 其中规模最大的两笔为：2024 年 4 月，银行授信，公司及子公司 2024 年度拟向银行等金融机构申请授信额度不超过人民币 10.00 亿元（含本数）；2021 年 12 月，可转债，发行 700.00 万张，每张面值 100.00 元。 股本方面，近三年披露股本变动 17 次（可转债转股 9 次、其他 3 次、股份回购 2 次、限售股份上市 1 次）；总股本由 2021 年末的 1.66 亿股变为 1.66 亿股（-0.29\%）。 分红方面，近三年现金分红 6 次、合计每股派现 1.28 元，最近一次实施在 2026 年 6 月。 近三年“再融资”类公告 46 条，涉及定向增发、可转债、募集资金使用。 综合看，公司融资结构上以债务滚动为主、股权融资为补充；2025 年筹资活动现金流净额 -1.60 亿元，融资节奏仍以维持存量债务周转为主，与前述经营与投资现金流的方向基本一致。

\pfl{未来潜在融资需求分析} 2026 年 6 月末货币资金 3.59 亿元，而短期债务 2.02 亿元（现金短债比 1.77 倍），2026 年上半年经营现金流净额 1.28 亿元。账面货币资金可以覆盖短期债务，缺口不体现在静态偿付层面。 按第 \ref{sec:financing-need} 节的三档划分，公司属于第三档“扩张型”：2026 年上半年 ROE 9.1\%（未年化，折合约 18.1\%）、6 月末资产负债率 31.8\%，资产负债表仍具备加杠杆空间；融资用途以产能扩张、品类并购与海外布局为主，单笔规模通常在 5---20 亿元量级，会被市场定价为成长而非补血。 公开披露的后续资金安排线索包括：股权融资线索：2026 年 5 月，公司 2024 年度向特定对象发行股票，已因股东会未通过延期而失效；债务与授信安排：2024 年 4 月，公司及子公司向金融机构申请综合授信额度不超过人民币 10.00 亿元。 资金需求还要叠加在建工程：期末在建工程 1.58 亿元、2026 年上半年购建固定资产等支出 2.35 亿元（未年化），这部分支出在周期底部通常会被主动放缓。 综合看，公司在静态偿付层面不存在缺口，外部融资更多是配合扩张节奏与并购整合的主动性安排，而非偿债压力驱动。




\subsubsection{金河生物（002688）}
\label{co:002688}

\pfl{公司基本情况} 内蒙古托克托县起家的饲用金霉素龙头，金霉素产销量约占全球一半，2012 年 7 月在深交所中小板上市，现形成兽用化药、兽用疫苗、环保污水处理与玉米深加工业务格局；2016 年、2021 年两次定增合计募资约 12.9 亿元，2025 年因中止非洲猪瘟疫苗研发全额计提吉林佑本商誉减值 1.75 亿元、归母净利润降至 2,755.67 万元，2026 年 8 月再以简易程序定增募资 3 亿元。 公司于 2012-07-13 在深交所主板首发上市，注册资本 8.37 亿元，总股本 7.72 亿股。 截至 2026 年 9 月 24 日收盘价 5.15 元，流通市值 38.7 亿元，近一年-23.4\%，流通市值在三级行业“动物保健Ⅲ”的 9 家公司中按由大到小排第\zhnumber{6}位，近一年涨跌幅低于全样本中位数（-10.1\%）13.3 个百分点。 按 2026 年半年报披露的行业构成，收入占比前三位为兽用化学药品 60.6\%；农产品加工业 22.3\%；兽用生物制品 8.8\%。

\pfl{历史股价} 以 2021 年 1 月为起点、2026 年 9 月为终点，股价由 5.05 元变为 5.15 元，累计 +2.0\%，区间最高 7.12 元（2021 年 6 月）、最低 3.63 元（2024 年 6 月）；当前价相当于区间最高价的 72.3\%，为区间最低价的 1.42 倍。 月度收益的年化波动率 30.2\%，低于全样本中位数 37.4\%，在三级行业“动物保健Ⅲ”的 9 家公司中波动率由高到低排第\zhnumber{9}位。 把股价阶段与公司事件对齐看，2024 年 2 月 至 2025 年 8 月股价上涨+94.5\%，同期（2023 年 12 月）公司收购子公司产生的商誉账面价值为 4.99 亿元。 整体上看，区间的几处大级别拐点都能在公司的资本运作、股权变动或风险事件上找到对应，股价波动有明显的公司层面驱动，而非单纯的行业 beta。

\begin{center}
\begin{minipage}{\textwidth}\centering
\begin{tikzpicture}
\begin{axis}[
  width=12.6cm, height=3.9cm, scale only axis,
  xmin=2020.922, xmax=2026.888, ymin=3.38, ymax=7.62,
  xtick={2021,2022,2023,2024,2025,2026},
  yearticks,
  yticklabel style={font=\scriptsize},
  ylabel={元/股}, ylabel style={font=\scriptsize},
  grid=major, grid style={LightGray, line width=0.3pt},
  axis line style={InkGray, line width=0.5pt},
  every axis plot/.append style={line width=0.9pt},
]
\addplot[AgriGreen] table[x=x,y=y]{data/clean/profiles/px_002688.dat};
\addplot[SoftRed, dashed, line width=0.7pt] table[x=x,y=y]{data/clean/profiles/pxzz_002688.dat};
\addplot[only marks, mark=*, mark size=1.1pt, SoftRed] table[x=x,y=y]{data/clean/profiles/pxzz_002688.dat};
\end{axis}
\end{tikzpicture}

\captionof{figure}[金河生物（002688）股价与周期骨架]{金河生物（002688）月度收盘价（元/股）与 ZigZag 周期骨架（阈值 25\%，圆点为周期转折点）}
\end{minipage}
\end{center}

\pfl{过去周期分析} 以 25\% 的 ZigZag 阈值划分，样本区间内股价形成 3 个主升段与 3 个主跌段。 最大主升段出现在 2024 年 2 月 至 2025 年 8 月，18 个月+94.5\%；最大主跌段为 2021 年 6 月 至 2022 年 4 月，10 个月-40.9\%。 最近一次转折出现在 2026 年 6 月（下跌段自 2025 年 8 月起，10 个月-34.9\%），当前处于该段的延续之中。 公司股价较申万农林牧渔指数提前约 23 个月见底，是板块中较早定价周期反转的公司之一。 动保的需求由存栏量而非价格决定（第 \ref{sec:vet} 节）：猪价下行阶段龙头养殖企业的逆势扩产反而推升了疫苗与兽药需求，这正是动保股价与猪价“脱钩”的原因，公司 2026 年上半年实现盈利、半年 ROE 5.7\%（未年化），好于全样本中位数（0.75\%），下行周期对其报表的冲击小于同业。 综合区间形态看，该股单段涨跌幅度偏大、以趋势段为主（最大主升 +94.5\%、最大主跌 -40.9\%），年化波动率 30.2\% 与全样本中位数 37.4\% 相比波动小于样本中位数。

\pfl{盈利情况} 2026 年上半年营业收入 14.48 亿元（同比 +4.1\%），归母净利润 1.31 亿元，同比 -5.1\%。 以年报为趋势对照，2023---2025 年营业收入由 21.74 亿元变为 28.78 亿元（两年年均复合增速 +15.1\%），归母净利润由 0.86 亿元变为 0.28 亿元。 按半年报口径，2026 年上半年的 ROE 为 5.7\%（未年化）、销售净利率 9.3\%、毛利率 37.8\%，ROE 在“动物保健Ⅲ”9 家公司中按数值由高到低排第\zhnumber{3}位。 按同一口径，三级行业“动物保健Ⅲ”9 家公司的半年 ROE 中位数为 2.95\%，该公司高于其中位数 2.79 个百分点。 从公司披露的原因看，业绩变动主要来自：2023 年 12 月，公司收购子公司产生的商誉账面价值为 4.99 亿元。 综合看，公司在报告期维持盈利（高于全样本半年 ROE 中位数 0.75\%），利润的可持续性取决于动物保健Ⅲ景气的延续与成本管控的执行。

\begin{center}\small\setlength{\tabcolsep}{3.4pt}
\captionof{table}[金河生物（002688）关键财务指标]{金河生物（002688）关键财务指标（合并报表口径；两个半年报期均未年化；现金短债比 $=$ 货币资金 $\div$ 短期债务）}
\begin{adjustbox}{max width=\textwidth}
\begin{tabular}{@{}lrrrrr@{}}
\toprule
指标 & 2023 年 & 2024 年 & 2025 年 & 2025 年半年报 & 2026 年半年报 \\
\midrule
营业收入（亿元） & 21.74 & 23.71 & 28.78 & 13.90 & 14.48 \\
归母净利润（亿元） & 0.86 & 1.00 & 0.28 & 1.38 & 1.31 \\
毛利率（\%） & 29.6 & 33.5 & 33.9 & 34.6 & 37.8 \\
ROE（\%） & 3.9 & 4.5 & 1.3 & 6.0 & 5.7 \\
资产负债率（\%） & 53.9 & 55.8 & 55.7 & 55.5 & 56.1 \\
经营活动现金流净额（亿元） & 1.99 & 4.98 & 4.69 & 1.34 & 2.54 \\
货币资金（亿元） & 6.08 & 6.36 & 6.18 & 6.25 & 7.00 \\
有息负债合计（亿元） & 23.00 & 23.19 & 23.20 & 24.15 & 23.64 \\
现金短债比（倍） & 0.34 & 0.48 & 0.47 & 0.41 & 0.51 \\
\bottomrule
\end{tabular}
\end{adjustbox}
\end{center}

\begin{center}
\begin{minipage}{\textwidth}\centering
\begin{tikzpicture}
\begin{groupplot}[
  group style={group size=2 by 1, horizontal sep=0.60cm},
  width=6.85cm, height=4.60cm,
  tick label style={font=\tiny},
  title style={font=\scriptsize\heiti},
  yticklabel style={font=\tiny},
  ylabel style={font=\tiny},
  grid=major, grid style={LightGray, line width=0.3pt},
  axis line style={InkGray, line width=0.5pt},
  enlarge x limits={abs=0.8},
  legend style={font=\scriptsize, at={(0.5,-0.20)}, anchor=north,
                legend columns=2, draw=none, fill=none, column sep=6pt},
]
\nextgroupplot[
  ybar, bar width=4pt,
  ymin=-2.88, ymax=32.23,
  xtick={0,1,2,3,4,5.7},
  xticklabels={2021,2022,2023,2024,2025,2026H1},
  ylabel={亿元},
]
\addplot[fill=AgriGreen!75, draw=AgriGreen, bar shift=-2.6pt] table[x=i, y=rev]{data/clean/profiles/fin_002688.dat};
\addplot[fill=SoftRed!60, draw=SoftRed, bar shift=2.6pt] table[x=i, y=ni]{data/clean/profiles/fin_002688.dat};
\addplot[fill=AgriGreen!30, draw=AgriGreen, pattern=north east lines, bar shift=-2.6pt] table[x=x, y=rev]{data/clean/profiles/finh_002688.dat};
\addplot[fill=SoftRed!20, draw=SoftRed, pattern=north east lines, bar shift=2.6pt] table[x=x, y=ni]{data/clean/profiles/finh_002688.dat};
\legend{营业收入（年/半年）, 归母净利润（年/半年）}
\nextgroupplot[
  ymin=-4.5, ymax=61.7,
  xtick={0,1,2,3,4,5.7},
  xticklabels={2021,2022,2023,2024,2025,2026H1},
  ylabel={\%},
]
\addplot[AgriGold, mark=*, mark size=1.3pt] table[x=i, y=roe]{data/clean/profiles/fin_002688.dat};
\addplot[OilBlue, mark=square*, mark size=1.3pt] table[x=i, y=debt]{data/clean/profiles/fin_002688.dat};
\addplot[only marks, AgriGold, mark=*, mark size=2.0pt] table[x=x, y=roe]{data/clean/profiles/finh_002688.dat};
\addplot[only marks, OilBlue, mark=square*, mark size=2.0pt] table[x=x, y=debt]{data/clean/profiles/finh_002688.dat};
\legend{ROE, 资产负债率}
\end{groupplot}
\end{tikzpicture}
\begin{tikzpicture}
\begin{axis}[
  name=finbot,
  width=13.75cm, height=3.10cm, scale only axis,
  ymin=-2.36, ymax=26.47,
  xtick={0,1,2,3,4,5.7},
  xticklabels={2021,2022,2023,2024,2025,2026H1},
  tick label style={font=\tiny},
  yticklabel style={font=\tiny},
  ylabel={亿元}, ylabel style={font=\tiny},
  ybar, bar width=4pt,
  grid=major, grid style={LightGray, line width=0.3pt},
  axis line style={InkGray, line width=0.5pt},
  enlarge x limits={abs=0.8},
  legend style={font=\scriptsize, at={(0.5,-0.26)}, anchor=north,
                legend columns=2, draw=none, fill=none, column sep=10pt},
]
\addplot[fill=AgriGreen!55, draw=AgriGreen, bar shift=-3.0pt] table[x=i, y=cash]{data/clean/profiles/fin_002688.dat};
\addplot[fill=SoftRed!45, draw=SoftRed, bar shift=3.0pt] table[x=i, y=ibd]{data/clean/profiles/fin_002688.dat};
\addplot[fill=AgriGreen!25, draw=AgriGreen, pattern=north east lines, bar shift=-3.0pt] table[x=x, y=cash]{data/clean/profiles/finh_002688.dat};
\addplot[fill=SoftRed!20, draw=SoftRed, pattern=north east lines, bar shift=3.0pt] table[x=x, y=ibd]{data/clean/profiles/finh_002688.dat};
\legend{货币资金（左轴）, 有息负债（左轴）}
\end{axis}
\begin{axis}[
  at={(finbot.south east)}, anchor=south east,
  width=13.75cm, height=3.10cm, scale only axis,
  axis y line*=right, axis x line*=none, xtick=\empty,
  ymin=-1.30, ymax=6.48,
  tick label style={font=\tiny},
  yticklabel style={font=\tiny}, scaled y ticks=false,
  legend style={font=\scriptsize, at={(0.5,-0.44)}, anchor=north,
                legend columns=1, draw=none, fill=none},
]
\addplot[OilBlue, mark=*, mark size=2.2pt, line width=1.3pt] table[x=i, y=ocf]{data/clean/profiles/fin_002688.dat};
\addplot[only marks, OilBlue, mark=*, mark size=3.2pt] table[x=x, y=ocf]{data/clean/profiles/finh_002688.dat};
\legend{经营活动现金流净额（右轴，亿元，蓝点）}
\end{axis}
\end{tikzpicture}

\captionof{figure}[金河生物（002688）近年主要财务指标]{金河生物（002688）近年主要财务指标（左上：营业收入与归母净利润；右上：ROE 与资产负债率；下：货币资金、有息负债为左轴柱状，经营活动现金流净额为其右轴的蓝点折线。
2021---2025 年为年报口径，2026H1 为 2026 年半年报/半年末，与其前一年报之间留出间隔、以斜纹或空心点区分；半年报数据未年化）}
\end{minipage}
\end{center}

\pfl{财务分析} 2026 年 6 月末资产负债率 56.1\%，较 2025 年末上升 0.4 个百分点（样本中位数 52.2\%，所处三级行业内按由低到高排第\zhnumber{8}位）。 2026 年 6 月末货币资金 7.00 亿元、短期债务（短期借款与一年内到期非流动负债）13.82 亿元、长期债务 9.82 亿元，有息负债合计 23.64 亿元，现金短债比 0.51 倍——覆盖偏紧。 净有息负债 16.64 亿元，相当于其 2026 年上半年营业收入的 115\%。 流动比率 1.18、速动比率 0.77，短期流动性无明显缺口。 2026 年上半年经营活动现金流净额 2.54 亿元，经营现金流与利润同向为正，内生资金可以支撑运营。 2026 年 6 月末在建工程 2.83 亿元，较上年末增加 0.87 亿元，仍有在建投入。 2026 年 6 月末商誉 2.78 亿元，占归母净资产的 11.2\%，是净资产中最脆弱的一块。 综合看，现金短债比 0.51 倍，短期资金安排偏紧；资产负债率高于样本中位数 3.9 个百分点；有息负债中短期债务占比更高，期限结构仍以滚动续作为主。 公司披露的风险集中在诉讼与资产冻结、股东与质押两类：一是 2026 年半年报，显示公司货币资金 2，819 元（母公司口径列示于受限资产表）因吉林佑本供应商起诉导致账户被冻结；二是 2026 年，一季报显示控股股东内蒙古金河控股有限公司持有 2.19 亿股（占 28.49\%），其中质押 8。

\pfl{重大战略转型} 公司的战略动作可归为以下几类：融资与资本运作（1 项）：2018 年 1 月，公司拟公开发行不超过 6.00 亿元可转债投建动物疫苗生产基地建设项目；并购与重组（1 项）：2023 年 5 月，控股子公司金河佑本以股权转让及增资方式取得吉林百思万可 60\% 股权。 公告口径上，近三年（2023 年 9 月---2026 年 9 月）公司披露重大重组与并购类公告 1 条、对外投资与转型类公告 5 条、回购与股权激励类公告 65 条，合计公告 581 条。 第一大行业仍是“兽用化学药品”，收入占比 60.6\%，与 2021 年年报相比没有方向性变化。 综合判断，报告期内公司有 6 条与战略相关的公告落地（重大重组与并购 1 条、对外投资与转型 5 条），资本运作与主业调整之间存在直接关联。

\begin{center}\small\setlength{\tabcolsep}{4pt}
\captionof{table}[金河生物（002688）历史融资与资本运作]{金河生物（002688）历史融资与资本运作（据公司公告与公开报道整理；金额为披露值，含首次公开发行、再融资、债券、银行借款与重整投资等）}
\begin{adjustbox}{max width=\textwidth}
\begin{tabular}{@{}p{1.45cm}p{2.55cm}p{4.05cm}p{4.95cm}@{}}
\toprule
时间 & 方式 & 规模 & 用途与说明 \\
\midrule
2012-07-05 & IPO（深圳证券交易所中小板） & 发行 2,\allowbreak{}723 万股，\allowbreak{}发行价 18.00 元/\allowbreak{}股，\allowbreak{}募集资金总额 4.90 亿元 & 年产 10，\allowbreak{}000 吨高效饲用金霉素扩建项目、\allowbreak{}金河生物研发中心建设项目；经证监发行字[2012]838 号文核准 \\
2016-02-03 & 非公开发行股票 & 发行 3，\allowbreak{}627.59 万股，\allowbreak{}发行价 13.05 元/\allowbreak{}股，\allowbreak{}认购资金总额 4.73 亿元 & 全部用于补充流动资金；；2016 年 1 月 21 日至 1 月 22 日完成验资 \\
2018-01-04 & 公开发行可转换公司债券（终止） & 拟募集资金总额（含发行费用）不超过人民币 6.00 亿元 & 动物疫苗生产基地建设项目（项目总投资 6.90 亿元 \\
2021-08-10 & 非公开发行股票（向 10 名认购对象） & 发行 1.45 亿股，\allowbreak{}发行价 5.65 元/\allowbreak{}股，\allowbreak{}募集资金总额 8.20 亿元 & 募集资金于 2021 年 7 月 5 日全部到位 \\
2023-05-24 & 现金收购与增资（并购） & 控股子公司金河佑本以 1.63 亿元受让吉林百思万可（后更名金河佑本（吉林）生物科技有限公司）34\% 股权 & 布局非洲猪瘟防控及疫苗产研领域；标的公司原股东为李忠祥、\allowbreak{}张年红、\allowbreak{}张维等 15 名自然人；2026 年 1 月 30 日公司决定中止非洲猪瘟冻干灭活疫苗研发和申报事项 \\
2026-01-30 & 控股子公司售后回租融资租赁（公司提供担保） & 公司按持股比例 85.6071\% 为金河佑本售后回租融资租赁业务提供担保 & 拓宽控股子公司金河佑本融资渠道、\allowbreak{}优化融资结构；经公司第六届董事会第三十六次会议（2026 年 1 月 30 日）审议通过并提交 2026 年第二次临时股东会审议；同日董事会还审议通过变更募集资金用途，\allowbreak{}颗粒剂自动化密闭生产线项目” \\
2026-08-27 & 以简易程序向特定对象发行股票（定增） & 发行 6，\allowbreak{}787.33 万股，\allowbreak{}发行价 4.42 元/\allowbreak{}股，\allowbreak{}募集资金总额 3.00 亿元 & 污水处理扩容及水资源再生利用提标扩建工程项目（拟投入 1.36 亿元）、\allowbreak{}新建产品库房建设项目（粮食储备库 \\
\bottomrule
\end{tabular}
\end{adjustbox}
\end{center}

\pfl{历史融投资情况} 从现金流量表看，2025 年公司取得借款收到的现金 14.17 亿元、偿还债务支付的现金 14.26 亿元、吸收权益性投资 0.02 亿元，筹资活动现金流净额 -1.24 亿元（2023 年为取得借款 20.33 亿元、偿还债务 14.33 亿元）。 2026 年上半年筹资活动现金流净额为 -0.24 亿元（取得借款 5.90 亿元、偿还债务 6.75 亿元，未年化），已转为净流出，处于净还债状态。 2025 年分配股利、利润或偿付利息支付的现金合计 1.48 亿元。 公开披露的融资记录共 7 笔，工具类型以 IPO、定向增发、可转债、融资租赁为主；金额与用途见下表。 其中规模最大的两笔为：2026 年 1 月，融资租赁，公司按持股比例 85.6071\% 为金河佑本售后回租融资租赁业务提供担保；2012 年 7 月，IPO，发行 2,723 万股，发行价 18.00 元/股，募集资金总额 4.90 亿元。 股本方面，近三年披露股本变动 7 次（股份回购 3 次、股权激励 2 次、其他 1 次）；总股本由 2021 年末的 7.80 亿股变为 7.72 亿股（-1.12\%）。 分红方面，近三年现金分红 4 次、合计每股派现 0.40 元，最近一次实施在 2025 年 12 月。 近三年“再融资”类公告 26 条，涉及定向增发、募集资金使用。 综合看，公司融资结构上以债务滚动为主、股权融资为补充；2025 年筹资活动现金流净额 -1.24 亿元，融资节奏仍以维持存量债务周转为主，与前述经营与投资现金流的方向基本一致。

\pfl{未来潜在融资需求分析} 2026 年 6 月末货币资金 7.00 亿元，而短期债务 13.82 亿元（现金短债比 0.51 倍），2026 年上半年经营现金流净额 2.54 亿元。按“短期债务减去账面货币资金”的滚动偿付口径，静态缺口约 6.82 亿元；上半年盈利或接近盈亏平衡，该缺口不随经营亏损进一步扩大。 按第 \ref{sec:financing-need} 节的三档划分，公司属于第三档“扩张型”：2026 年上半年 ROE 5.7\%（未年化，折合约 11.5\%）、6 月末资产负债率 56.1\%，资产负债表仍具备加杠杆空间；融资用途以产能扩张、品类并购与海外布局为主，单笔规模通常在 5---20 亿元量级，会被市场定价为成长而非补血。 公开披露的后续资金安排线索包括：股权融资线索：2026 年 8 月，公司完成 2025 年度以简易程序向特定对象发行股票；债务与授信安排：2026 年 1 月，控股子公司金河佑本开展售后回租融资租赁业务、公司按 85.6071\% 持股比例提供不超过 1.21 亿元担保并由控股股东全额担保。 资金需求还要叠加在建工程：期末在建工程 2.83 亿元、2026 年上半年购建固定资产等支出 1.44 亿元（未年化），这部分支出在周期底部通常会被主动放缓。 综合看，公司的外部资金需求以补充营运资金为主、项目投入为辅，滚动偿付口径下的静态缺口约 6.82 亿元，融资节奏将受制于银行续贷意愿与股权融资窗口。




\subsubsection{普莱柯（603566）}
\label{co:603566}

\pfl{公司基本情况} 河南洛阳的综合性兽用疫苗与兽药企业，2002 年成立、2015 年 5 月在上交所上市，产品覆盖猪用、禽用疫苗及兽用化药；2022 年 9 月完成 8.98 亿元非公开发行投向兽用灭活疫苗产能，2025 年归母净利润 1.85 亿元、同比增长 99.33\%，但 2026 年上半年受下游生猪价格低位、动保行业竞争加剧及生物制品增值税率上调影响，营收 4.91 亿元、同比下降 12.18\%，归母净利润同比下降 33.45\%。 公司于 2015-05-18 在上交所首发上市，注册资本 3.46 亿元，总股本 3.46 亿股。 截至 2026 年 9 月 24 日收盘价 10.19 元，流通市值 35.3 亿元，近一年-25.3\%，流通市值在三级行业“动物保健Ⅲ”的 9 家公司中按由大到小排第\zhnumber{7}位，近一年涨跌幅低于全样本中位数（-10.1\%）15.2 个百分点。 按 2025 年年报披露的行业构成，收入占比前三位为兽用药品 94.0\%；技术许可或转让 2.5\%；功能性保健品等(宠物用) 2.3\%。

\pfl{历史股价} 本报告样本区间（2021 年 1 月 至 2026 年 9 月）内，股价由 22.31 元变为 10.19 元，累计 -54.3\%，区间最高 31.89 元（2023 年 2 月）、最低 9.79 元（2026 年 6 月）；当前价相当于区间最高价的 32.0\%，已接近区间最低价（1.04 倍）。 月度收益的年化波动率 34.9\%，低于全样本中位数 37.4\%，在三级行业“动物保健Ⅲ”的 9 家公司中波动率由高到低排第\zhnumber{7}位。 把股价阶段与公司事件对齐看，2023 年 2 月 至 2026 年 6 月股价下跌-69.3\%，同期（2026 年 6 月）非公开发行募集资金净额 8.86 亿元截至 2026 年 6 月末累计投入 6.51 亿元（进度 73.44\%）、变更用途募资总额 33.07 万元。 把这些阶段与公司事件对齐后可以看到，几轮大涨大跌多由公司自身的资本运作与风险事件触发，行业景气只解释了其中一部分。

\begin{center}
\begin{minipage}{\textwidth}\centering
\begin{tikzpicture}
\begin{axis}[
  width=12.6cm, height=3.9cm, scale only axis,
  xmin=2020.922, xmax=2026.888, ymin=9.10, ymax=34.12,
  xtick={2021,2022,2023,2024,2025,2026},
  yearticks,
  yticklabel style={font=\scriptsize},
  ylabel={元/股}, ylabel style={font=\scriptsize},
  grid=major, grid style={LightGray, line width=0.3pt},
  axis line style={InkGray, line width=0.5pt},
  every axis plot/.append style={line width=0.9pt},
]
\addplot[AgriGreen] table[x=x,y=y]{data/clean/profiles/px_603566.dat};
\addplot[SoftRed, dashed, line width=0.7pt] table[x=x,y=y]{data/clean/profiles/pxzz_603566.dat};
\addplot[only marks, mark=*, mark size=1.1pt, SoftRed] table[x=x,y=y]{data/clean/profiles/pxzz_603566.dat};
\end{axis}
\end{tikzpicture}

\captionof{figure}[普莱柯（603566）股价与周期骨架]{普莱柯（603566）月度收盘价（元/股）与 ZigZag 周期骨架（阈值 25\%，圆点为周期转折点）}
\end{minipage}
\end{center}

\pfl{过去周期分析} 以 25\% 的 ZigZag 阈值划分，样本区间内股价形成 2 个主升段与 3 个主跌段。 最大主跌段为 2023 年 2 月 至 2026 年 6 月，40 个月-69.3\%；最大主升段出现在 2021 年 10 月 至 2022 年 2 月，4 个月+58.7\%。 最近一次转折出现在 2026 年 6 月（下跌段自 2023 年 2 月起，40 个月-69.3\%），当前处于该段的延续之中。 公司股价与申万农林牧渔指数几乎同步见底，个股并未走出独立行情。 动保的需求由存栏量而非价格决定（第 \ref{sec:vet} 节）：猪价下行阶段龙头养殖企业的逆势扩产反而推升了疫苗与兽药需求，这正是动保股价与猪价“脱钩”的原因，公司 2026 年上半年实现盈利、半年 ROE 3.0\%（未年化），好于全样本中位数（0.75\%），下行周期对其报表的冲击小于同业。 综合区间形态看，该股单段涨跌幅度偏大、以趋势段为主（最大主升 +58.7\%、最大主跌 -69.3\%），年化波动率 34.9\% 与全样本中位数 37.4\% 相比波动与样本中位数接近。

\pfl{盈利情况} 2026 年上半年营业收入 4.91 亿元（同比 -12.2\%），归母净利润 0.77 亿元，同比 -33.4\%。 以年报为趋势对照，2023---2025 年营业收入由 12.53 亿元变为 10.90 亿元（两年年均复合增速 -6.7\%），归母净利润由 1.75 亿元变为 1.85 亿元。 按半年报口径，2026 年上半年的 ROE 为 3.0\%（未年化）、销售净利率 13.0\%、毛利率 55.8\%，ROE 在“动物保健Ⅲ”9 家公司中按数值由高到低排第\zhnumber{5}位。 按同一口径，三级行业“动物保健Ⅲ”9 家公司的半年 ROE 中位数为 2.95\%，该公司低于其中位数 0.00 个百分点。 从公司披露的原因看，业绩变动主要来自：2026 年上半年，归属于上市公司股东的净利润 7，732.42 万元、同比下降 33.45\%；2026 年半年报，主要系下游生猪出栏价格长期处于低位、动保行业竞争加剧、猪用疫苗销量和价格下降，同时生物制品增值税适用税率由 3\% 上调至 13\%。 面对这一局面，公司披露的举措包括：2024 年 9 月，公司存在因企业合并形成的商誉（收购普莱柯（南京）形成商誉资产组。 综合看，公司在报告期维持盈利（高于全样本半年 ROE 中位数 0.75\%），利润的可持续性取决于动物保健Ⅲ景气的延续与成本管控的执行。

\begin{center}\small\setlength{\tabcolsep}{3.4pt}
\captionof{table}[普莱柯（603566）关键财务指标]{普莱柯（603566）关键财务指标（合并报表口径；两个半年报期均未年化；现金短债比 $=$ 货币资金 $\div$ 短期债务）}
\begin{adjustbox}{max width=\textwidth}
\begin{tabular}{@{}lrrrrr@{}}
\toprule
指标 & 2023 年 & 2024 年 & 2025 年 & 2025 年半年报 & 2026 年半年报 \\
\midrule
营业收入（亿元） & 12.53 & 10.43 & 10.90 & 5.59 & 4.91 \\
归母净利润（亿元） & 1.75 & 0.93 & 1.85 & 1.16 & 0.77 \\
毛利率（\%） & 61.0 & 61.1 & 59.9 & 59.8 & 55.8 \\
ROE（\%） & 6.3 & 3.5 & 7.1 & 4.4 & 3.0 \\
资产负债率（\%） & 16.5 & 15.2 & 15.4 & 14.8 & 15.2 \\
经营活动现金流净额（亿元） & 2.83 & 2.69 & 3.56 & 1.95 & 0.78 \\
货币资金（亿元） & 3.17 & 2.05 & 1.37 & 1.20 & 0.86 \\
有息负债合计（亿元） & 0.02 & 0.01 & 0.27 & 0.01 & 0.26 \\
现金短债比（倍） & 318.06 & 189.62 & 19.77 & 110.46 & 11.86 \\
\bottomrule
\end{tabular}
\end{adjustbox}
\end{center}

\begin{center}
\begin{minipage}{\textwidth}\centering
\begin{tikzpicture}
\begin{groupplot}[
  group style={group size=2 by 1, horizontal sep=0.60cm},
  width=6.85cm, height=4.60cm,
  tick label style={font=\tiny},
  title style={font=\scriptsize\heiti},
  yticklabel style={font=\tiny},
  ylabel style={font=\tiny},
  grid=major, grid style={LightGray, line width=0.3pt},
  axis line style={InkGray, line width=0.5pt},
  enlarge x limits={abs=0.8},
  legend style={font=\scriptsize, at={(0.5,-0.20)}, anchor=north,
                legend columns=2, draw=none, fill=none, column sep=6pt},
]
\nextgroupplot[
  ybar, bar width=4pt,
  ymin=-1.25, ymax=14.03,
  xtick={0,1,2,3,4,5.7},
  xticklabels={2021,2022,2023,2024,2025,2026H1},
  ylabel={亿元},
]
\addplot[fill=AgriGreen!75, draw=AgriGreen, bar shift=-2.6pt] table[x=i, y=rev]{data/clean/profiles/fin_603566.dat};
\addplot[fill=SoftRed!60, draw=SoftRed, bar shift=2.6pt] table[x=i, y=ni]{data/clean/profiles/fin_603566.dat};
\addplot[fill=AgriGreen!30, draw=AgriGreen, pattern=north east lines, bar shift=-2.6pt] table[x=x, y=rev]{data/clean/profiles/finh_603566.dat};
\addplot[fill=SoftRed!20, draw=SoftRed, pattern=north east lines, bar shift=2.6pt] table[x=x, y=ni]{data/clean/profiles/finh_603566.dat};
\legend{营业收入（年/半年）, 归母净利润（年/半年）}
\nextgroupplot[
  ymin=-1.4, ymax=19.1,
  xtick={0,1,2,3,4,5.7},
  xticklabels={2021,2022,2023,2024,2025,2026H1},
  ylabel={\%},
]
\addplot[AgriGold, mark=*, mark size=1.3pt] table[x=i, y=roe]{data/clean/profiles/fin_603566.dat};
\addplot[OilBlue, mark=square*, mark size=1.3pt] table[x=i, y=debt]{data/clean/profiles/fin_603566.dat};
\addplot[only marks, AgriGold, mark=*, mark size=2.0pt] table[x=x, y=roe]{data/clean/profiles/finh_603566.dat};
\addplot[only marks, OilBlue, mark=square*, mark size=2.0pt] table[x=x, y=debt]{data/clean/profiles/finh_603566.dat};
\legend{ROE, 资产负债率}
\end{groupplot}
\end{tikzpicture}
\begin{tikzpicture}
\begin{axis}[
  name=finbot,
  width=13.75cm, height=3.10cm, scale only axis,
  ymin=-0.32, ymax=3.55,
  xtick={0,1,2,3,4,5.7},
  xticklabels={2021,2022,2023,2024,2025,2026H1},
  tick label style={font=\tiny},
  yticklabel style={font=\tiny},
  ylabel={亿元}, ylabel style={font=\tiny},
  ybar, bar width=4pt,
  grid=major, grid style={LightGray, line width=0.3pt},
  axis line style={InkGray, line width=0.5pt},
  enlarge x limits={abs=0.8},
  legend style={font=\scriptsize, at={(0.5,-0.26)}, anchor=north,
                legend columns=2, draw=none, fill=none, column sep=10pt},
]
\addplot[fill=AgriGreen!55, draw=AgriGreen, bar shift=-3.0pt] table[x=i, y=cash]{data/clean/profiles/fin_603566.dat};
\addplot[fill=SoftRed!45, draw=SoftRed, bar shift=3.0pt] table[x=i, y=ibd]{data/clean/profiles/fin_603566.dat};
\addplot[fill=AgriGreen!25, draw=AgriGreen, pattern=north east lines, bar shift=-3.0pt] table[x=x, y=cash]{data/clean/profiles/finh_603566.dat};
\addplot[fill=SoftRed!20, draw=SoftRed, pattern=north east lines, bar shift=3.0pt] table[x=x, y=ibd]{data/clean/profiles/finh_603566.dat};
\legend{货币资金（左轴）, 有息负债（左轴）}
\end{axis}
\begin{axis}[
  at={(finbot.south east)}, anchor=south east,
  width=13.75cm, height=3.10cm, scale only axis,
  axis y line*=right, axis x line*=none, xtick=\empty,
  ymin=-0.93, ymax=4.63,
  tick label style={font=\tiny},
  yticklabel style={font=\tiny}, scaled y ticks=false,
  legend style={font=\scriptsize, at={(0.5,-0.44)}, anchor=north,
                legend columns=1, draw=none, fill=none},
]
\addplot[OilBlue, mark=*, mark size=2.2pt, line width=1.3pt] table[x=i, y=ocf]{data/clean/profiles/fin_603566.dat};
\addplot[only marks, OilBlue, mark=*, mark size=3.2pt] table[x=x, y=ocf]{data/clean/profiles/finh_603566.dat};
\legend{经营活动现金流净额（右轴，亿元，蓝点）}
\end{axis}
\end{tikzpicture}

\captionof{figure}[普莱柯（603566）近年主要财务指标]{普莱柯（603566）近年主要财务指标（左上：营业收入与归母净利润；右上：ROE 与资产负债率；下：货币资金、有息负债为左轴柱状，经营活动现金流净额为其右轴的蓝点折线。
2021---2025 年为年报口径，2026H1 为 2026 年半年报/半年末，与其前一年报之间留出间隔、以斜纹或空心点区分；半年报数据未年化）}
\end{minipage}
\end{center}

\pfl{财务分析} 2026 年 6 月末资产负债率 15.2\%，较 2025 年末下降 0.2 个百分点（样本中位数 52.2\%，所处三级行业内按由低到高排第\zhnumber{3}位）。 2026 年 6 月末货币资金 0.86 亿元、短期债务（短期借款与一年内到期非流动负债）0.07 亿元、长期债务 0.19 亿元，有息负债合计 0.26 亿元，现金短债比 11.86 倍——覆盖充分。 净有息负债 -0.60 亿元，相当于其 2026 年上半年营业收入的 -12\%。 流动比率 3.46、速动比率 2.86，短期指标尚可。 2026 年上半年经营活动现金流净额 0.78 亿元，经营现金流与利润同向为正，内生资金可以支撑运营。 综合看，现金短债比 11.86 倍，短期偿付压力可控；资产负债率低于样本中位数 37.0 个百分点；有息负债中长期债务占比更高，期限结构对短债滚动的压力较小。

\pfl{重大战略转型} 近三年公司的战略动作集中在：融资与资本运作（1 项）：2020 年 6 月，公司以非公开发行募集资金投向兽用灭活疫苗生产项目。 公告口径上，近三年（2023 年 9 月---2026 年 9 月）公司披露重大重组与并购类公告 2 条、对外投资与转型类公告 0 条、回购与股权激励类公告 20 条，合计公告 306 条。 第一大行业仍是“兽用药品”，收入占比 94.0\%，与 2021 年年报相比没有方向性变化。 综合判断，报告期内公司有 2 条与战略相关的公告落地（重大重组与并购 2 条、对外投资与转型 0 条），资本运作与主业调整之间存在直接关联。

\begin{center}\small\setlength{\tabcolsep}{4pt}
\captionof{table}[普莱柯（603566）历史融资与资本运作]{普莱柯（603566）历史融资与资本运作（据公司公告与公开报道整理；金额为披露值，含首次公开发行、再融资、债券、银行借款与重整投资等）}
\begin{adjustbox}{max width=\textwidth}
\begin{tabular}{@{}p{1.45cm}p{2.55cm}p{4.05cm}p{4.95cm}@{}}
\toprule
时间 & 方式 & 规模 & 用途与说明 \\
\midrule
2015-05-18 & IPO（上海证券交易所主板） & 发行 4,\allowbreak{}000 万股，\allowbreak{}发行价 15.52 元/\allowbreak{}股；2015 年募集资金总额 5.60 亿元 & 动物疫苗产业化建设项目（新增 6 条生产线用于 17 个产品生产）、\allowbreak{}技术研发中心建设等项目；采用网下向投资者询价配售与网上按市值申购定价发行相结合的方式，\allowbreak{}发行后总股本不超 1.6 亿股；2015 年 8 月 20 日披露的募集资金使用情况显示募集资金总额 5.60 亿元 \\
2020-06-20 & 非公开发行股票（获核准批文） & 拟发行股票募集资金总额不超过 8.98 亿元 & 兽用灭活疫苗生产项目等；其中全资子公司普莱柯（南京）拟投资 4.00 亿元建设兽用灭活疫苗生产项目（一期） \\
2022-09-05 & 非公开发行股票（向特定对象） & 发行 3，\allowbreak{}142.06 万股，\allowbreak{}发行价 28.58 元/\allowbreak{}股，\allowbreak{}募集资金总额 8.98 亿元 & 兽用灭活疫苗生产项目等（截至 2026 年 6 月末累计投入 6.51 亿元、\allowbreak{}投入进度 73.44\%）；200 股；投资者于 2022 年 9 月 2 日 17:00 前足额缴款 \\
\bottomrule
\end{tabular}
\end{adjustbox}
\end{center}

\pfl{历史融投资情况} 从现金流量表看，2025 年公司，筹资活动现金流净额 -2.07 亿元（2023 年为偿还债务 0.07 亿元）。 2026 年上半年筹资活动现金流净额为 -1.37 亿元（取得借款 0.25 亿元、偿还债务 0.25 亿元，未年化），已转为净流出，处于净还债状态。 2025 年分配股利、利润或偿付利息支付的现金合计 2.05 亿元。 公开披露的融资记录共 3 笔，工具类型以 IPO、定向增发为主；金额与用途见下表。 其中规模最大的两笔为：2015 年 5 月，IPO，发行 4,000 万股，发行价 15.52 元/股，2015 年募集资金总额 5.60 亿元；2020 年 6 月，定向增发，拟发行股票募集资金总额不超过 8.98 亿元。 股本方面，近三年披露股本变动 8 次（限售股份上市 2 次、其他 2 次、股份回购 1 次）；总股本由 2021 年末的 3.21 亿股变为 3.46 亿股（+7.64\%）。 分红方面，近三年现金分红 6 次、合计每股派现 2.05 元，最近一次实施在 2025 年 12 月。 近三年“再融资”类公告 4 条，涉及定向增发。 综合看，公司融资结构上股权与债务融资并举；2025 年筹资活动现金流净额 -2.07 亿元，融资节奏仍以维持存量债务周转为主，与前述经营与投资现金流的方向基本一致。

\pfl{未来潜在融资需求分析} 2026 年 6 月末货币资金 0.86 亿元，而短期债务 0.07 亿元（现金短债比 11.86 倍），2026 年上半年经营现金流净额 0.78 亿元。账面货币资金可以覆盖短期债务，缺口不体现在静态偿付层面。 公司 2026 年 6 月末资产负债率 15.2\%、2026 年上半年 ROE 3.0\%（未年化），资产端与负债端都没有明显压力，融资需求以相机抉择的扩张型用途为主。 公开披露的后续资金安排线索包括：股权融资线索：2026 年 6 月，非公开发行募集资金净额 8.86 亿元截至 2026 年 6 月末累计投入 6.51 亿元（进度 73.44\%）、变更用途募资总额 33.07 万元。 静态偿付不存在缺口，后续融资更可能服务于产能或并购安排，而不是填补流动性窟窿。




\subsubsection{申联生物（688098）}
\label{co:688098}

\pfl{公司基本情况} 2018 年设立至今，申联生物是专注口蹄疫疫苗的科创板动保企业（上海），2019 年 10 月 28 日上市、为首家科创板动保公司，主要产品为口蹄疫合成肽疫苗与口蹄疫灭活疫苗；上市后业绩持续承压并处于破发状态，2025 年营业收入 2.88 亿元、归母净利润 -2,049.50 万元（亏损收窄约 54\%），2026 年 2 月以 23,744.1789 万元收购扬州世之源控股权、开始构建“人用药品+动物保健”双主业。 公司于 2019-10-28 在上交所科创板首发上市，注册资本 4.11 亿元，总股本 4.11 亿股。 截至 2026 年 9 月 24 日收盘价 8.37 元，流通市值 34.4 亿元，近一年-22.5\%，流通市值在三级行业“动物保健Ⅲ”的 9 家公司中按由大到小排第\zhnumber{8}位，近一年涨跌幅低于全样本中位数（-10.1\%）12.4 个百分点。 按 2025 年年报披露的行业构成，收入占比前三位为兽用生物制品 99.3\%；其他 0.5\%；其他(补充) 0.2\%。

\pfl{历史股价} 以 2021 年 1 月为起点、2026 年 9 月为终点，股价由 14.49 元变为 8.37 元，累计 -42.2\%，区间最高 14.61 元（2021 年 2 月）、最低 3.92 元（2024 年 8 月）；当前价相当于区间最高价的 57.3\%，为区间最低价的 2.14 倍。 月度收益的年化波动率 38.3\%，高于全样本中位数 37.4\%，在三级行业“动物保健Ⅲ”的 9 家公司中波动率由高到低排第\zhnumber{4}位。

\begin{center}
\begin{minipage}{\textwidth}\centering
\begin{tikzpicture}
\begin{axis}[
  width=12.6cm, height=3.9cm, scale only axis,
  xmin=2020.922, xmax=2026.888, ymin=3.65, ymax=15.63,
  xtick={2021,2022,2023,2024,2025,2026},
  yearticks,
  yticklabel style={font=\scriptsize},
  ylabel={元/股}, ylabel style={font=\scriptsize},
  grid=major, grid style={LightGray, line width=0.3pt},
  axis line style={InkGray, line width=0.5pt},
  every axis plot/.append style={line width=0.9pt},
]
\addplot[AgriGreen] table[x=x,y=y]{data/clean/profiles/px_688098.dat};
\addplot[SoftRed, dashed, line width=0.7pt] table[x=x,y=y]{data/clean/profiles/pxzz_688098.dat};
\addplot[only marks, mark=*, mark size=1.1pt, SoftRed] table[x=x,y=y]{data/clean/profiles/pxzz_688098.dat};
\end{axis}
\end{tikzpicture}

\captionof{figure}[申联生物（688098）股价与周期骨架]{申联生物（688098）月度收盘价（元/股）与 ZigZag 周期骨架（阈值 25\%，圆点为周期转折点）}
\end{minipage}
\end{center}

\pfl{过去周期分析} 以 25\% 的 ZigZag 阈值划分，样本区间内股价形成 1 个主升段与 2 个主跌段。 最大主升段出现在 2024 年 8 月 至 2025 年 8 月，12 个月+215.8\%；最大主跌段为 2021 年 1 月 至 2024 年 8 月，43 个月-72.9\%。 最近一次转折出现在 2026 年 5 月（下跌段自 2025 年 8 月起，9 个月-37.2\%），当前处于该段的延续之中。 公司股价较申万农林牧渔指数提前约 21 个月见底，是板块中较早定价周期反转的公司之一。 动保的需求由存栏量而非价格决定（第 \ref{sec:vet} 节）：猪价下行阶段龙头养殖企业的逆势扩产反而推升了疫苗与兽药需求，这正是动保股价与猪价“脱钩”的原因，公司 2026 年 6 月末资产负债率 9.8\%、半年 ROE -1.4\%（未年化），在本轮下行中属于随行业同步调整的一类。 综合区间形态看，该股单段涨跌幅度偏大、以趋势段为主（最大主升 +215.8\%、最大主跌 -72.9\%），年化波动率 38.3\% 与全样本中位数 37.4\% 相比波动与样本中位数接近。

\pfl{盈利情况} 2026 年上半年营业收入 1.12 亿元（同比 -8.1\%），归母净利润 -0.19 亿元，亏损同比扩大 0.06 亿元。 以年报为趋势对照，2023---2025 年营业收入由 3.01 亿元变为 2.88 亿元（两年年均复合增速 -2.2\%），归母净利润由 0.32 亿元变为 -0.20 亿元。 按半年报口径，2026 年上半年的 ROE 为 -1.4\%（未年化）、销售净利率 -18.2\%、毛利率 57.4\%，ROE 在“动物保健Ⅲ”9 家公司中按数值由高到低排第\zhnumber{9}位。 按同一口径，三级行业“动物保健Ⅲ”9 家公司的半年 ROE 中位数为 2.95\%，该公司低于其中位数 4.34 个百分点。 从公司披露的原因看，业绩变动主要来自：2025 年，归属于上市公司股东的净利润 -2，049.50 万元，原因是母公司存在未弥补亏损。 综合看，公司仍处亏损区间、亏损同比扩大，半年 ROE 在三级行业内按由高到低排第\zhnumber{9}位，单位盈利能力的修复依赖行业景气与自身成本端的改善，报表弹性主要体现为价格与销量的方向性变化。

\begin{center}\small\setlength{\tabcolsep}{3.4pt}
\captionof{table}[申联生物（688098）关键财务指标]{申联生物（688098）关键财务指标（合并报表口径；两个半年报期均未年化；现金短债比 $=$ 货币资金 $\div$ 短期债务）}
\begin{adjustbox}{max width=\textwidth}
\begin{tabular}{@{}lrrrrr@{}}
\toprule
指标 & 2023 年 & 2024 年 & 2025 年 & 2025 年半年报 & 2026 年半年报 \\
\midrule
营业收入（亿元） & 3.01 & 3.03 & 2.88 & 1.22 & 1.12 \\
归母净利润（亿元） & 0.32 & -0.45 & -0.20 & -0.13 & -0.19 \\
毛利率（\%） & 70.8 & 60.2 & 56.8 & 58.8 & 57.4 \\
ROE（\%） & 2.1 & -3.1 & -1.5 & -0.9 & -1.4 \\
资产负债率（\%） & 5.5 & 8.3 & 9.3 & 8.3 & 9.8 \\
经营活动现金流净额（亿元） & 0.04 & 0.89 & 0.33 & -0.39 & -0.65 \\
货币资金（亿元） & 0.45 & 1.46 & 1.22 & 0.53 & 0.65 \\
有息负债合计（亿元） & 0.01 & 0.15 & 0.14 & 0.15 & 0.53 \\
现金短债比（倍） & 134.73 & 10.09 & 9.20 & 3.80 & 1.23 \\
\bottomrule
\end{tabular}
\end{adjustbox}
\end{center}

\begin{center}
\begin{minipage}{\textwidth}\centering
\begin{tikzpicture}
\begin{groupplot}[
  group style={group size=2 by 1, horizontal sep=0.60cm},
  width=6.85cm, height=4.60cm,
  tick label style={font=\tiny},
  title style={font=\scriptsize\heiti},
  yticklabel style={font=\tiny},
  ylabel style={font=\tiny},
  grid=major, grid style={LightGray, line width=0.3pt},
  axis line style={InkGray, line width=0.5pt},
  enlarge x limits={abs=0.8},
  legend style={font=\scriptsize, at={(0.5,-0.20)}, anchor=north,
                legend columns=2, draw=none, fill=none, column sep=6pt},
]
\nextgroupplot[
  ybar, bar width=4pt,
  ymin=-0.85, ymax=4.07,
  xtick={0,1,2,3,4,5.7},
  xticklabels={2021,2022,2023,2024,2025,2026H1},
  ylabel={亿元},
]
\addplot[fill=AgriGreen!75, draw=AgriGreen, bar shift=-2.6pt] table[x=i, y=rev]{data/clean/profiles/fin_688098.dat};
\addplot[fill=SoftRed!60, draw=SoftRed, bar shift=2.6pt] table[x=i, y=ni]{data/clean/profiles/fin_688098.dat};
\addplot[fill=AgriGreen!30, draw=AgriGreen, pattern=north east lines, bar shift=-2.6pt] table[x=x, y=rev]{data/clean/profiles/finh_688098.dat};
\addplot[fill=SoftRed!20, draw=SoftRed, pattern=north east lines, bar shift=2.6pt] table[x=x, y=ni]{data/clean/profiles/finh_688098.dat};
\legend{营业收入（年/半年）, 归母净利润（年/半年）}
\nextgroupplot[
  ymin=-4.1, ymax=11.1,
  xtick={0,1,2,3,4,5.7},
  xticklabels={2021,2022,2023,2024,2025,2026H1},
  ylabel={\%},
]
\addplot[AgriGold, mark=*, mark size=1.3pt] table[x=i, y=roe]{data/clean/profiles/fin_688098.dat};
\addplot[OilBlue, mark=square*, mark size=1.3pt] table[x=i, y=debt]{data/clean/profiles/fin_688098.dat};
\addplot[only marks, AgriGold, mark=*, mark size=2.0pt] table[x=x, y=roe]{data/clean/profiles/finh_688098.dat};
\addplot[only marks, OilBlue, mark=square*, mark size=2.0pt] table[x=x, y=debt]{data/clean/profiles/finh_688098.dat};
\legend{ROE, 资产负债率}
\end{groupplot}
\end{tikzpicture}
\begin{tikzpicture}
\begin{axis}[
  name=finbot,
  width=13.75cm, height=3.10cm, scale only axis,
  ymin=-0.15, ymax=1.64,
  xtick={0,1,2,3,4,5.7},
  xticklabels={2021,2022,2023,2024,2025,2026H1},
  tick label style={font=\tiny},
  yticklabel style={font=\tiny},
  ylabel={亿元}, ylabel style={font=\tiny},
  ybar, bar width=4pt,
  grid=major, grid style={LightGray, line width=0.3pt},
  axis line style={InkGray, line width=0.5pt},
  enlarge x limits={abs=0.8},
  legend style={font=\scriptsize, at={(0.5,-0.26)}, anchor=north,
                legend columns=2, draw=none, fill=none, column sep=10pt},
]
\addplot[fill=AgriGreen!55, draw=AgriGreen, bar shift=-3.0pt] table[x=i, y=cash]{data/clean/profiles/fin_688098.dat};
\addplot[fill=SoftRed!45, draw=SoftRed, bar shift=3.0pt] table[x=i, y=ibd]{data/clean/profiles/fin_688098.dat};
\addplot[fill=AgriGreen!25, draw=AgriGreen, pattern=north east lines, bar shift=-3.0pt] table[x=x, y=cash]{data/clean/profiles/finh_688098.dat};
\addplot[fill=SoftRed!20, draw=SoftRed, pattern=north east lines, bar shift=3.0pt] table[x=x, y=ibd]{data/clean/profiles/finh_688098.dat};
\legend{货币资金（左轴）, 有息负债（左轴）}
\end{axis}
\begin{axis}[
  at={(finbot.south east)}, anchor=south east,
  width=13.75cm, height=3.10cm, scale only axis,
  axis y line*=right, axis x line*=none, xtick=\empty,
  ymin=-1.05, ymax=1.37,
  tick label style={font=\tiny},
  yticklabel style={font=\tiny}, scaled y ticks=false,
  legend style={font=\scriptsize, at={(0.5,-0.44)}, anchor=north,
                legend columns=1, draw=none, fill=none},
]
\addplot[OilBlue, mark=*, mark size=2.2pt, line width=1.3pt] table[x=i, y=ocf]{data/clean/profiles/fin_688098.dat};
\addplot[only marks, OilBlue, mark=*, mark size=3.2pt] table[x=x, y=ocf]{data/clean/profiles/finh_688098.dat};
\legend{经营活动现金流净额（右轴，亿元，蓝点）}
\end{axis}
\end{tikzpicture}

\captionof{figure}[申联生物（688098）近年主要财务指标]{申联生物（688098）近年主要财务指标（左上：营业收入与归母净利润；右上：ROE 与资产负债率；下：货币资金、有息负债为左轴柱状，经营活动现金流净额为其右轴的蓝点折线。
2021---2025 年为年报口径，2026H1 为 2026 年半年报/半年末，与其前一年报之间留出间隔、以斜纹或空心点区分；半年报数据未年化）}
\end{minipage}
\end{center}

\pfl{财务分析} 2026 年 6 月末资产负债率 9.8\%，较 2025 年末上升 0.5 个百分点（样本中位数 52.2\%，所处三级行业内按由低到高排第\zhnumber{1}位）。 2026 年 6 月末货币资金 0.65 亿元、短期债务（短期借款与一年内到期非流动负债）0.53 亿元、长期债务 0.01 亿元，有息负债合计 0.53 亿元，现金短债比 1.23 倍——覆盖充分。 净有息负债 -0.11 亿元，相当于其 2026 年上半年营业收入的 -10\%。 流动比率 3.54、速动比率 2.72，短期流动性无明显缺口。 2026 年上半年经营活动现金流净额 -0.65 亿元，经营现金流与利润同时为负，现金消耗较快。 综合看，现金短债比 1.23 倍，短期偿付压力可控；资产负债率低于样本中位数 42.4 个百分点；有息负债中短期债务占比更高，期限结构仍以滚动续作为主。

\pfl{重大战略转型} 近三年公司的战略动作集中在：产能与项目（1 项）：2025 年 12 月，公司 2025 年疫苗销售头份数同比增长 2.42\%、亏损幅度收窄约 54\%；并购与重组（1 项）：2026 年 2 月，公司收购扬州世之源控股权，投资方式为股权受让+增资并引入员工持股平台。 公告口径上，近三年（2023 年 9 月---2026 年 9 月）公司披露重大重组与并购类公告 2 条、对外投资与转型类公告 3 条、回购与股权激励类公告 22 条，合计公告 243 条。 第一大行业仍是“兽用生物制品”，收入占比 99.3\%，与 2021 年年报相比没有方向性变化。 综合判断，报告期内公司有 5 条与战略相关的公告落地（重大重组与并购 2 条、对外投资与转型 3 条），资本运作与主业调整之间存在直接关联。

\begin{center}\small\setlength{\tabcolsep}{4pt}
\captionof{table}[申联生物（688098）历史融资与资本运作]{申联生物（688098）历史融资与资本运作（据公司公告与公开报道整理；金额为披露值，含首次公开发行、再融资、债券、银行借款与重整投资等）}
\begin{adjustbox}{max width=\textwidth}
\begin{tabular}{@{}p{1.45cm}p{2.55cm}p{4.05cm}p{4.95cm}@{}}
\toprule
时间 & 方式 & 规模 & 用途与说明 \\
\midrule
2019-10-16 & IPO（上海证券交易所科创板） & 发行 5，\allowbreak{}000 万股（不低于发行后总股本的 10.00\% & 悬浮培养口蹄疫灭活疫苗项目（招股说明书原拟募集资金 4.50 亿元）；发行市盈率（按发行后总股本）43.87 倍 \\
2023-08-30 & 首发募集资金使用（募投项目） & 募集资金 4.40 亿元、\allowbreak{}净额 400，\allowbreak{}175，\allowbreak{}0 元量级 & 增加并扩大公司口蹄疫疫苗种类及生产规模；公司预先以自筹资金开展募投项目 \\
2026-02-25 & 收购联营公司控股权并开展新业务（关联交易，含银行并购贷款安排） & 公司拟使用 2.37 亿元的自有资金及自筹资金（包括银行并购贷款） & 加速推进公司创新药战略布局 \\
\bottomrule
\end{tabular}
\end{adjustbox}
\end{center}

\pfl{历史融投资情况} 从现金流量表看，2025 年公司偿还债务支付的现金 0.01 亿元，筹资活动现金流净额 -0.04 亿元。 2026 年上半年筹资活动现金流净额为 0.55 亿元（取得借款 0.40 亿元、偿还债务 0.01 亿元，未年化），仍处于净融入状态。 2025 年分配股利、利润或偿付利息支付的现金合计 0.00 亿元。 公开披露的融资记录共 3 笔，工具类型以 IPO 为主；金额与用途见下表。 其中规模最大的两笔为：2026 年 2 月，收购联营公司控股权并，公司拟使用 2.37 亿元的自有资金及自筹资金（包括银行并购贷款）；2019 年 10 月，IPO，发行 5，000 万股（不低于发行后总股本的 10.00\%。 股本方面，近三年披露股本变动 3 次；总股本自 2021 年末以来未发生变化（4.11 亿股）。 分红方面，近三年现金分红 2 次、合计每股派现 0.11 元，最近一次实施在 2023 年 12 月。 近三年“再融资”类公告 3 条，涉及定向增发。 综合看，公司融资结构上股权与债务融资并举；2025 年筹资活动现金流净额 -0.04 亿元，融资节奏仍以维持存量债务周转为主，与前述经营与投资现金流的方向基本一致。

\pfl{未来潜在融资需求分析} 2026 年 6 月末货币资金 0.65 亿元，而短期债务 0.53 亿元（现金短债比 1.23 倍），2026 年上半年经营现金流净额 -0.65 亿元。账面货币资金可以覆盖短期债务，缺口不体现在静态偿付层面。 公司 2026 年上半年归母净利润为负（-0.19 亿元），但 6 月末资产负债率 9.8\% 尚未进入第一档（亏损且负债率 $>$65\%）的区间，当前融资需求以补充流动资金、支撑低谷期经营性支出为主。 公开披露的后续资金安排线索包括：债务与授信安排：2026 年 2 月，公司明确收购世之源控股权的 2.37 亿元资金来源于自有资金及自筹资金，其中包含银行并购贷款。 静态偿付不存在缺口，后续融资更可能服务于产能或并购安排，而不是填补流动性窟窿。




\subsubsection{绿康生化（002868）}
\label{co:002868}

\pfl{公司基本情况} 福建南平浦城起家的兽药（饲料添加剂、生物发酵）企业，2017 年 5 月深交所上市；2022—2023 年跨界收购光伏胶膜企业江西纬科并投建海宁胶膜项目，2022—2024 年连续三年扣非净利润为负、2024 年末经审计净资产为 -2,453.61 万元，2025 年 4 月 30 日起被实施退市风险警示（*ST 绿康）；2025 年 4 月原控股股东向福建纵腾网络转让 29.99\% 股份、实际控制人变更为王钻，公司以 0 元剥离三家光伏胶膜子公司，2026 年 5 月 13 日撤销退市风险警示及其他风险警示，2026 年上半年归母净利润 731.60 万元、扭亏为盈。 公司于 2017-05-03 在深交所主板首发上市，注册资本 1.55 亿元，总股本 1.55 亿股。 截至 2026 年 9 月 24 日收盘价 22.42 元，流通市值 24.1 亿元，近一年-19.9\%，流通市值在三级行业“动物保健Ⅲ”的 9 家公司中按由大到小排第\zhnumber{9}位，近一年涨跌幅低于全样本中位数（-10.1\%）9.8 个百分点。 按 2026 年半年报披露的行业构成，收入占比前三位为动保系列 91.4\%；其他 5.0\%；植保系列 3.7\%。 与 2021 年年报相比，第一大行业口径由“兽药”（67.9\%）变为“动保系列”（91.4\%）；两期披露的分类口径有调整，占比差异中同时包含分类因素。

\pfl{历史股价} 自 2021 年 1 月至 2026 年 9 月，股价由 11.13 元变为 22.42 元，累计 +101.4\%，区间最高 61.22 元（2023 年 1 月）、最低 9.29 元（2021 年 10 月）；当前价相当于区间最高价的 36.6\%，为区间最低价的 2.41 倍。 月度收益的年化波动率 57.4\%，高于全样本中位数 37.4\%，在三级行业“动物保健Ⅲ”的 9 家公司中波动率由高到低排第\zhnumber{1}位。 把股价阶段与公司事件对齐看，2021 年 1 月 至 2023 年 1 月股价上涨+450.0\%，同期（2023 年 1 月）董事会审议通过以 9，500.00 万元受让江西纬科 100\% 股权；2025 年 3 月 至 2026 年 2 月股价上涨+189.2\%，同期（2025 年）公司完成剥离胶膜业务、回归兽药业务。 从时间对应关系看，公司层面的资本运作与股权、风险事件解释了区间内多数急涨急跌，与行业指数的同步性相对有限。

\begin{center}
\begin{minipage}{\textwidth}\centering
\begin{tikzpicture}
\begin{axis}[
  width=12.6cm, height=3.9cm, scale only axis,
  xmin=2020.922, xmax=2026.888, ymin=8.64, ymax=65.51,
  xtick={2021,2022,2023,2024,2025,2026},
  yearticks,
  yticklabel style={font=\scriptsize},
  ylabel={元/股}, ylabel style={font=\scriptsize},
  grid=major, grid style={LightGray, line width=0.3pt},
  axis line style={InkGray, line width=0.5pt},
  every axis plot/.append style={line width=0.9pt},
]
\addplot[AgriGreen] table[x=x,y=y]{data/clean/profiles/px_002868.dat};
\addplot[SoftRed, dashed, line width=0.7pt] table[x=x,y=y]{data/clean/profiles/pxzz_002868.dat};
\addplot[only marks, mark=*, mark size=1.1pt, SoftRed] table[x=x,y=y]{data/clean/profiles/pxzz_002868.dat};
\end{axis}
\end{tikzpicture}

\captionof{figure}[绿康生化（002868）股价与周期骨架]{绿康生化（002868）月度收盘价（元/股）与 ZigZag 周期骨架（阈值 25\%，圆点为周期转折点）}
\end{minipage}
\end{center}

\pfl{过去周期分析} 以 25\% 的 ZigZag 阈值划分，样本区间内股价形成 3 个主升段与 3 个主跌段。 最大主升段出现在 2021 年 1 月 至 2023 年 1 月，24 个月+450.0\%；最大主跌段为 2023 年 1 月 至 2024 年 8 月，19 个月-79.9\%。 最近一次转折出现在 2026 年 7 月（下跌段自 2026 年 2 月起，5 个月-49.9\%），当前处于该段的延续之中。 公司股价较申万农林牧渔指数提前约 55 个月见底，是板块中较早定价周期反转的公司之一。 动保的需求由存栏量而非价格决定（第 \ref{sec:vet} 节）：猪价下行阶段龙头养殖企业的逆势扩产反而推升了疫苗与兽药需求，这正是动保股价与猪价“脱钩”的原因，公司 2026 年 6 月末资产负债率 90.1\%，高于全样本中位数 52.2\%，其股价因此更多反映资金面与信用风险，而不是价格弹性本身。 综合区间形态看，该股单段涨跌幅度偏大、以趋势段为主（最大主升 +450.0\%、最大主跌 -79.9\%），年化波动率 57.4\% 与全样本中位数 37.4\% 相比波动明显大于样本中位数。

\pfl{盈利情况} 2026 年上半年营业收入 2.48 亿元（同比 -15.7\%），归母净利润 0.07 亿元，实现扭亏。 以年报为趋势对照，2023---2025 年营业收入由 5.07 亿元变为 5.17 亿元（两年年均复合增速 +0.9\%），归母净利润由 -2.22 亿元变为 -1.39 亿元。 按半年报口径，2026 年上半年的 ROE 为 10.9\%（未年化）、销售净利率 2.8\%、毛利率 22.9\%，ROE 在“动物保健Ⅲ”9 家公司中按数值由高到低排第\zhnumber{1}位。 按同一口径，三级行业“动物保健Ⅲ”9 家公司的半年 ROE 中位数为 2.95\%，该公司高于其中位数 7.96 个百分点。 从公司披露的原因看，业绩变动主要来自：2025 年，公司完成剥离胶膜业务、回归兽药业务；2025 年 12 月，盐霉素、莫能菌素等抗球虫药产能逐步释放，实现新旧产能系统性转换与升级。 面对这一局面，公司披露的举措包括：2025 年 11 月，公司公告以 0 元价格完成出售三家光伏胶膜子公司。 综合看，公司在报告期维持盈利（高于全样本半年 ROE 中位数 0.75\%），利润的可持续性取决于动物保健Ⅲ景气的延续与成本管控的执行。

\begin{center}\small\setlength{\tabcolsep}{3.4pt}
\captionof{table}[绿康生化（002868）关键财务指标]{绿康生化（002868）关键财务指标（合并报表口径；两个半年报期均未年化；现金短债比 $=$ 货币资金 $\div$ 短期债务）}
\begin{adjustbox}{max width=\textwidth}
\begin{tabular}{@{}lrrrrr@{}}
\toprule
指标 & 2023 年 & 2024 年 & 2025 年 & 2025 年半年报 & 2026 年半年报 \\
\midrule
营业收入（亿元） & 5.07 & 6.49 & 5.17 & 2.94 & 2.48 \\
归母净利润（亿元） & -2.22 & -4.45 & -1.39 & -0.59 & 0.07 \\
毛利率（\%） & -9.0 & -8.4 & 6.7 & 5.9 & 22.9 \\
ROE（\%） & -45.1 & 0.0 & 170.5 & 0.0 & 10.9 \\
资产负债率（\%） & 78.9 & 101.6 & 91.3 & 105.8 & 90.1 \\
经营活动现金流净额（亿元） & -1.87 & 1.41 & -0.25 & 0.75 & 0.45 \\
货币资金（亿元） & 0.84 & 0.25 & 0.56 & 0.16 & 0.51 \\
有息负债合计（亿元） & 9.63 & 9.16 & 3.60 & 8.66 & 3.19 \\
现金短债比（倍） & 0.19 & 0.03 & 0.17 & 0.02 & 0.17 \\
\bottomrule
\end{tabular}
\end{adjustbox}
\end{center}

\begin{center}
\begin{minipage}{\textwidth}\centering
\begin{tikzpicture}
\begin{groupplot}[
  group style={group size=2 by 1, horizontal sep=0.60cm},
  width=6.85cm, height=4.60cm,
  tick label style={font=\tiny},
  title style={font=\scriptsize\heiti},
  yticklabel style={font=\tiny},
  ylabel style={font=\tiny},
  grid=major, grid style={LightGray, line width=0.3pt},
  axis line style={InkGray, line width=0.5pt},
  enlarge x limits={abs=0.8},
  legend style={font=\scriptsize, at={(0.5,-0.20)}, anchor=north,
                legend columns=2, draw=none, fill=none, column sep=6pt},
]
\nextgroupplot[
  ybar, bar width=4pt,
  ymin=-5.54, ymax=7.81,
  xtick={0,1,2,3,4,5.7},
  xticklabels={2021,2022,2023,2024,2025,2026H1},
  ylabel={亿元},
]
\addplot[fill=AgriGreen!75, draw=AgriGreen, bar shift=-2.6pt] table[x=i, y=rev]{data/clean/profiles/fin_002868.dat};
\addplot[fill=SoftRed!60, draw=SoftRed, bar shift=2.6pt] table[x=i, y=ni]{data/clean/profiles/fin_002868.dat};
\addplot[fill=AgriGreen!30, draw=AgriGreen, pattern=north east lines, bar shift=-2.6pt] table[x=x, y=rev]{data/clean/profiles/finh_002868.dat};
\addplot[fill=SoftRed!20, draw=SoftRed, pattern=north east lines, bar shift=2.6pt] table[x=x, y=ni]{data/clean/profiles/finh_002868.dat};
\legend{营业收入（年/半年）, 归母净利润（年/半年）}
\nextgroupplot[
  ymin=-62.4, ymax=192.0,
  xtick={0,1,2,3,4,5.7},
  xticklabels={2021,2022,2023,2024,2025,2026H1},
  ylabel={\%},
]
\addplot[AgriGold, mark=*, mark size=1.3pt] table[x=i, y=roe]{data/clean/profiles/fin_002868.dat};
\addplot[OilBlue, mark=square*, mark size=1.3pt] table[x=i, y=debt]{data/clean/profiles/fin_002868.dat};
\addplot[only marks, AgriGold, mark=*, mark size=2.0pt] table[x=x, y=roe]{data/clean/profiles/finh_002868.dat};
\addplot[only marks, OilBlue, mark=square*, mark size=2.0pt] table[x=x, y=debt]{data/clean/profiles/finh_002868.dat};
\legend{ROE, 资产负债率}
\end{groupplot}
\end{tikzpicture}
\begin{tikzpicture}
\begin{axis}[
  name=finbot,
  width=13.75cm, height=3.10cm, scale only axis,
  ymin=-0.96, ymax=10.78,
  xtick={0,1,2,3,4,5.7},
  xticklabels={2021,2022,2023,2024,2025,2026H1},
  tick label style={font=\tiny},
  yticklabel style={font=\tiny},
  ylabel={亿元}, ylabel style={font=\tiny},
  ybar, bar width=4pt,
  grid=major, grid style={LightGray, line width=0.3pt},
  axis line style={InkGray, line width=0.5pt},
  enlarge x limits={abs=0.8},
  legend style={font=\scriptsize, at={(0.5,-0.26)}, anchor=north,
                legend columns=2, draw=none, fill=none, column sep=10pt},
]
\addplot[fill=AgriGreen!55, draw=AgriGreen, bar shift=-3.0pt] table[x=i, y=cash]{data/clean/profiles/fin_002868.dat};
\addplot[fill=SoftRed!45, draw=SoftRed, bar shift=3.0pt] table[x=i, y=ibd]{data/clean/profiles/fin_002868.dat};
\addplot[fill=AgriGreen!25, draw=AgriGreen, pattern=north east lines, bar shift=-3.0pt] table[x=x, y=cash]{data/clean/profiles/finh_002868.dat};
\addplot[fill=SoftRed!20, draw=SoftRed, pattern=north east lines, bar shift=3.0pt] table[x=x, y=ibd]{data/clean/profiles/finh_002868.dat};
\legend{货币资金（左轴）, 有息负债（左轴）}
\end{axis}
\begin{axis}[
  at={(finbot.south east)}, anchor=south east,
  width=13.75cm, height=3.10cm, scale only axis,
  axis y line*=right, axis x line*=none, xtick=\empty,
  ymin=-2.72, ymax=2.40,
  tick label style={font=\tiny},
  yticklabel style={font=\tiny}, scaled y ticks=false,
  legend style={font=\scriptsize, at={(0.5,-0.44)}, anchor=north,
                legend columns=1, draw=none, fill=none},
]
\addplot[OilBlue, mark=*, mark size=2.2pt, line width=1.3pt] table[x=i, y=ocf]{data/clean/profiles/fin_002868.dat};
\addplot[only marks, OilBlue, mark=*, mark size=3.2pt] table[x=x, y=ocf]{data/clean/profiles/finh_002868.dat};
\legend{经营活动现金流净额（右轴，亿元，蓝点）}
\end{axis}
\end{tikzpicture}

\captionof{figure}[绿康生化（002868）近年主要财务指标]{绿康生化（002868）近年主要财务指标（左上：营业收入与归母净利润；右上：ROE 与资产负债率；下：货币资金、有息负债为左轴柱状，经营活动现金流净额为其右轴的蓝点折线。
2021---2025 年为年报口径，2026H1 为 2026 年半年报/半年末，与其前一年报之间留出间隔、以斜纹或空心点区分；半年报数据未年化）}
\end{minipage}
\end{center}

\pfl{财务分析} 2026 年 6 月末资产负债率 90.1\%，较 2025 年末下降 1.3 个百分点（样本中位数 52.2\%，所处三级行业内按由低到高排第\zhnumber{9}位）。 2026 年 6 月末货币资金 0.51 亿元、短期债务（短期借款与一年内到期非流动负债）2.94 亿元、长期债务 0.25 亿元，有息负债合计 3.19 亿元，现金短债比 0.17 倍——缺口明显。 净有息负债 2.67 亿元，相当于其 2026 年上半年营业收入的 108\%。 流动比率 0.37、速动比率 0.20，短期偿付压力较大。 2026 年上半年经营活动现金流净额 0.45 亿元，经营现金流与利润同向为正，内生资金可以支撑运营。 综合看，账面货币资金对有息负债与短期债务的覆盖不足，滚动偿付对新增融资的依赖度较高；资产负债率高于样本中位数 37.9 个百分点；有息负债中短期债务占比更高，期限结构仍以滚动续作为主。 公司披露的退市与持续经营风险为：2026 年 3 月，公司股票可能被终止上市的第三次风险提示公告；2025 年 4 月，公司股票自 该日起被实施退市风险警示、其他风险警示。

\pfl{重大战略转型} 近三年公司的战略动作集中在：并购与重组（5 项）：2026 年 5 月，公司股票撤销退市风险警示及其他风险警示；2025 年 9 月，公司深交所问询函回复；产能与项目（2 项）：2023 年 1 月，董事会审议通过以 9，500.00 万元受让江西纬科 100\% 股权。 公告口径上，近三年（2023 年 9 月---2026 年 9 月）公司披露重大重组与并购类公告 23 条、对外投资与转型类公告 5 条、回购与股权激励类公告 11 条，合计公告 642 条。 主营结构的口径变化已经落到报表上：2021 年年报的第一大行业为“兽药”（67.9\%），2026 年半年报为“动保系列”（91.4\%）；公司在这两期的分部划分方式有过调整。 综合判断，报告期内公司有 28 条与战略相关的公告落地（重大重组与并购 23 条、对外投资与转型 5 条），资本运作与主业调整之间存在直接关联。

\begin{center}\small\setlength{\tabcolsep}{4pt}
\captionof{table}[绿康生化（002868）历史融资与资本运作]{绿康生化（002868）历史融资与资本运作（据公司公告与公开报道整理；金额为披露值，含首次公开发行、再融资、债券、银行借款与重整投资等）}
\begin{adjustbox}{max width=\textwidth}
\begin{tabular}{@{}p{1.45cm}p{2.55cm}p{4.05cm}p{4.95cm}@{}}
\toprule
时间 & 方式 & 规模 & 用途与说明 \\
\midrule
2017-05-03 & IPO（深圳证券交易所） & 实际发行 3，\allowbreak{}000.00 万股 & 份有限公司，\allowbreak{}发行市盈率（按发行后总股本）20.50 倍 \\
2023-01-28 & 现金收购股权（跨界收购光伏胶膜企业） & 作价 9，\allowbreak{}溢价率约 632\%） & 由兽药跨界切入光伏胶膜赛道；经第四届董事会第二十一次（临时）会议审议通过；标的后更名为绿康（玉山）胶膜材料有限公司；公司另投资成立绿康（海宁）胶膜材料有限公司，\allowbreak{}计划在海宁投资 60.00 亿元建设年产 8 亿平方米光伏胶膜项目 \\
2023-06-22 & 向特定对象发行 A 股股票（定增预案，未实施） & 发行对象为不超过 35 名特定投资者，\allowbreak{}发行价格为不低于定价基准日前 20 个交易日公司股票交易均价的 80\% & 报告期无募集资金使用情况）；公司于 2023 年 6 月 22 日在巨潮资讯网披露绿康生化 2023 年度向特定对象发行 A 股股票预案；媒体 2025 年报道称公司定增推进不顺 \\
2024-06-27 & 出售全资子公司股权（关联交易） & 以 2，\allowbreak{}730.00 万元的交易价格向关联法人福建浦潭热能有限公司出售全资子公司武汉绿康生化科技有限公司 100\% 股权 & 剥离资产、\allowbreak{}回笼资金；经公司第五届董事会第二十六次会议审议 \\
2024-07-22 & 以简易程序向特定对象发行股票（定增预案，未实施） & 拟募集资金总额不超过 8，\allowbreak{}000.00 万元（符合简易程序募集资金不超过 3.00 亿元且不超过最近一年末净资产 20\% 的规定） & 投向绿康（海宁）胶膜材料有限公司年产 3.2 亿平方米光伏胶膜项目（项目投资总额 7.54 亿元 \\
2025-04-24 & 控制权转让（股份协议转让 29.99\%） & 康怡投资、\allowbreak{}义睿投资、\allowbreak{}长鑫贰号、\allowbreak{}皓赢投资向福建纵腾网络有限公司转让公司 4，\allowbreak{}660.84 万股普通股（占公司总股本的 29.99\%） & 引入产业投资人、\allowbreak{}化解经营与退市风险；2025 年 11 月 21 日该等股份完成过户 \\
2025-11 & 控股股东借款 & 向控股股东福建纵腾网络有限公司借入 1.50 亿元 & 用于公司日常生产经营；借款期限自借款协议签署日起计算；同期公司完成 1.26 亿元银行借款的到期续贷 \\
2025-11-26 & 重大资产出售（0 元剥离光伏胶膜资产） & 以 0 元价格向江西饶信新能材料有限公司出售绿康（玉山）胶膜材料有限公司、\allowbreak{}绿康（海宁）胶膜材料有限公司、\allowbreak{}绿康新能（上海）进出口贸易有限公司共三家标的公司 100\% 股权；以 2024 年 12 月 31 日为评估基准日 & 退出光伏胶膜赛道、\allowbreak{}剥离亏损资产；本次交易构成关联交易；2025 年 9 月 23 日公司深交所问询函回复 \\
2025-12 & 土地收储（政府补偿） & 公司及子公司福建绿安生物农药有限公司与浦城县土地储备中心签订浦城县国有土地使用权收购合同，\allowbreak{}收储价款人民币 1.87 亿元；截至问询函回复日土地收储补偿第一期款项 3 & 盘活土地资产、\allowbreak{}支持厂区搬迁与产能升级；收储价款将按协议分期结算 \\
\bottomrule
\end{tabular}
\end{adjustbox}
\end{center}

\pfl{历史融投资情况} 从现金流量表看，2025 年公司取得借款收到的现金 1.72 亿元、偿还债务支付的现金 5.07 亿元，筹资活动现金流净额 1.95 亿元（2023 年为取得借款 6.69 亿元、偿还债务 2.14 亿元）。 2026 年上半年筹资活动现金流净额为 -0.54 亿元（取得借款 1.73 亿元、偿还债务 2.04 亿元，未年化），已转为净流出，处于净还债状态。 2025 年分配股利、利润或偿付利息支付的现金合计 0.43 亿元。 公开披露的融资记录共 9 笔，工具类型以 IPO、定向增发、银行借款为主；金额与用途见下表。 其中规模最大的两笔为：2025 年 11 月，银行借款，向控股股东福建纵腾网络有限公司借入 1.50 亿元；2024 年 7 月，定向增发，拟募集资金总额不超过 8，000.00 万元（符合简易程序募集资金不超过 3.00 亿元且不超过最近一年末净资产 20\% 的规定）。 股本方面，近三年披露股本变动 3 次；总股本自 2021 年末以来未发生变化（1.55 亿股）。 分红方面，近三年未实施现金分红，在全部样本中属于分红能力偏弱的一端。 近三年“再融资”类公告 8 条，涉及定向增发、募集资金使用。 综合看，公司融资结构上股权与债务融资并举；2025 年筹资活动现金流净额 1.95 亿元，年度内仍为净融入，与前述经营与投资现金流的方向基本一致。

\pfl{未来潜在融资需求分析} 2026 年 6 月末货币资金 0.51 亿元，而短期债务 2.94 亿元（现金短债比 0.17 倍），2026 年上半年经营现金流净额 0.45 亿元。按“短期债务减去账面货币资金”的滚动偿付口径，静态缺口约 2.42 亿元；上半年盈利或接近盈亏平衡，该缺口不随经营亏损进一步扩大。 按第 \ref{sec:financing-need} 节的三档划分，公司属于第二档“保壳型”：近三年重大重组与并购类公告 23 条，2026 年上半年归母净利润 0.07 亿元，融资需求更多体现为“资产置换 + 维持上市地位”，而不是扩产。 公开披露的后续资金安排线索包括：股权融资线索：2024 年 7 月，公司 2024 年度以简易程序向特定对象发行股票预案；债务与授信安排：2025 年 11 月，同期到期银行借款续贷 1.26 亿元；2026 年 4 月，公司计划 2026 年向浦城农行申请流动资金 5，030.00 万元贷款续期；资本开支安排：2025 年 12 月，公司与浦城县土地储备中心签订浦城县国有土地使用权收购合同。 综合看，公司的外部资金需求以营运资金周转与在建项目投入为主，滚动偿付口径下的静态缺口约 2.42 亿元，融资节奏将受制于银行续贷意愿与股权融资窗口。




